% concatenated from main.tex
\documentclass[11pt,reqno]{amsart}
\usepackage{comment}
\usepackage{amssymb}
\usepackage{tikz-cd}
\usepackage{tikz}
\usepackage{amsthm}
\usepackage{graphicx}
\usepackage{caption}
\usepackage{subcaption}
\usepackage[backref]{hyperref}
\usepackage{cleveref} % Required for inserting images
\usepackage{booktabs}
\usepackage{float}

%uncomment to disable todonotes
\usepackage[disable]{todonotes}
%comment \iffalse \fi to enable todonotes
\iffalse
%\usepackage{todonotes}
%make todonotes work with amsart
\makeatletter
\providecommand\@dotsep{5}
\renewcommand{\listoftodos}[1][\@todonotes@todolistname]{%
	\@starttoc{tdo}{#1}}
\makeatother
\fi

\DeclareMathOperator{\rk}{rk}
\DeclareMathOperator{\Aut}{Aut}
\DeclareMathOperator{\Sym}{Sym}
\DeclareMathOperator{\cond}{cond}
\DeclareMathOperator{\End}{End}
\DeclareMathOperator{\Gal}{Gal}
\DeclareMathOperator{\HH}{H}
\DeclareMathOperator{\lcm}{lcm}
\DeclareMathOperator{\disc}{disc}
\DeclareMathOperator{\PSL}{PSL}
\DeclareMathOperator{\GL}{GL}
\DeclareMathOperator{\PGL}{PGL}
\DeclareMathOperator{\id}{id}
\DeclareMathOperator{\Id}{Id}
\DeclareMathOperator{\SL}{SL}
\DeclareMathOperator{\U}{U}
\DeclareMathOperator{\Cl}{Cl}
\DeclareMathOperator{\vol}{vol}
\DeclareMathOperator{\Hom}{Hom}
\DeclareMathOperator{\covol}{covol}
\DeclareMathOperator{\mat}{\mathcal{M}}
\DeclareMathOperator{\Frob}{Frob}
\DeclareMathOperator{\Cls}{Cls}
\DeclareMathOperator{\aCls}{Cls^{*1}}
\DeclareMathOperator{\trd}{trd}
\DeclareMathOperator{\nrd}{nrd}
\DeclareMathOperator{\proj}{pr}
\DeclareMathOperator{\Tr}{Tr}
\DeclareMathOperator{\Nm}{Nm}
\DeclareMathOperator{\Ver}{Ver}
\DeclareMathOperator{\tr}{tr}
\DeclareMathOperator{\diag}{diag}
\DeclareMathOperator{\dlog}{dlog}
\DeclareMathOperator{\Span}{span}
\DeclareMathOperator{\Sturm}{Sturm}
\DeclareMathOperator{\ord}{ord}
\DeclareMathOperator{\ind}{ind}
\DeclareMathOperator{\new}{\mathrm{new}}
\DeclareMathOperator{\old}{\mathrm{old}}
\DeclareMathOperator{\impart}{Im}
\DeclareMathOperator{\ModSym}{ModSym}
\DeclareMathOperator{\Div}{Div}
\DeclareMathOperator{\ddiv}{div}
\DeclareMathOperator{\Frac}{Frac}
\DeclareMathOperator{\red}{red}
\DeclareMathOperator{\gon}{gon}
\DeclareMathOperator{\img}{im}
\DeclareMathOperator{\Spec}{Spec}
\DeclareMathOperator{\Ell}{Ell}

\newtheorem{theorem}{Theorem}[section]
\newtheorem*{theorem*}{Theorem}
\newtheorem{lemma}[theorem]{Lemma}
\newtheorem{conjecture}[theorem]{Conjecture}
\newtheorem{proposition}[theorem]{Proposition}
\newtheorem*{proposition*}{Proposition}
\newtheorem{corollary}[theorem]{Corollary}
\newtheorem*{corollary*}{Corollary}
\newtheorem{claim}[theorem]{Claim}
\newtheorem{question}[theorem]{Question}
\theoremstyle{definition}
\newtheorem{remark}[theorem]{Remark}
\newtheorem{definition}[theorem]{Definition}
\newtheorem{example}[theorem]{Example}
\newtheorem{exercise}[theorem]{Exercise}
\numberwithin{equation}{section}

\newcommand{\modcrvgt}{\,_{>\,}\!}
\newcommand{\Q}{\mathbb{Q}}
\newcommand{\Z}{\mathbb{Z}}
\newcommand{\N}{\mathbb{N}}
\newcommand{\Zhat}{\widehat{\mathbb{Z}}}
\newcommand{\R}{\mathbb{R}}
\newcommand{\C}{\mathbb{C}}
\newcommand{\F}{\mathbb{F}}
\newcommand{\PP}{\mathbb{P}}
\newcommand{\TT}{\mathcal{T}}
\newcommand{\Qbar}{\Q^{\textup{al}}}
\newcommand{\calP}{\mathcal{P}}
\newcommand{\Hamil}{\mathbb{H}}
\newcommand{\frakm}{\mathfrak{m}}
\newcommand{\frakl}{\mathfrak{l}}
\newcommand{\eps}{\varepsilon}
\newcommand{\fp}{\mathfrak{p}}
\newcommand{\order}{\mathcal{O}}
\newcommand{\OO}{\mathcal{O}}
\newcommand{\calO}{\order}
\newcommand{\lat}{\mathrm{Lat}}
\newcommand{\tor}{\mathrm{tors}}
\newcommand{\hlat}{\lat^{*}}
%\newcommand{\hlat}{\mathrm{HLat}}
\newcommand{\hlato}{\lat^{*1}}
%\newcommand{\hlato}{\mathrm{HLat}^{1}}
\newcommand{\lambdanew}{\lambda_{\mathrm{new}}}
\def\torz#1{\mathbb Z /#1 \mathbb Z}
\def\tg#1#2{\mathbb Z/#1\mathbb Z \times \mathbb Z /#2 \mathbb Z}
\newcommand{\diamondop}[1]{\langle #1 \rangle} % diamond operator
\newcommand{\set}[1]{\left\lbrace #1 \right\rbrace}
\newcommand{\multiset}[1]{\llbrace #1 \rrbrace}
\newcommand{\llbrace}{\lbrace \!\!\!\: \lbrace}
\newcommand{\rrbrace}{\rbrace \!\!\!\: \rbrace}

\newcommand{\lmfdbec}[3]{\href{https://www.lmfdb.org/EllipticCurve/Q/#1/#2/#3}{#1.#2#3}}
\newcommand{\lmfdbeciso}[2]{\href{https://www.lmfdb.org/EllipticCurve/Q/#1/#2}{#1.#2}}
\newcommand{\lmfdbecnf}[4]{\href{https://www.lmfdb.org/EllipticCurve/#1/#2/#3/#4}{#1-#2-#3#4}}
\newcommand{\lmfdbecnfiso}[3]{\href{https://www.lmfdb.org/EllipticCurve/#1/#2/#3}{#1-#2-#3}}
\newcommand{\lmfdcharorbit}[2]{\href{https://www.lmfdb.org/Character/Dirichlet/#1/#2}{#1.#2}}
\newcommand{\lmfdbnewform}[4]{\href{https://www.lmfdb.org/ModularForm/GL2/Q/holomorphic/#1/#2/#3/#4}{#1.#2.#3.#4}}
\newcommand{\githubbare}[1]{\href{https://github.com/nt-lib/quartic-torsion/blob/main/#1}{\path{#1}}}
\newcommand{\github}[2]{[\githubbare{#1}, #2]}


\newcommand{\medit}[1]{{\color{blue} \sf  #1}}
\newcommand{\fedit}[1]{{\color{purple} \sf  #1}}
\newcommand{\maarten}[1]{{\color{blue} \sf $\clubsuit\clubsuit\clubsuit$ Maarten: [#1]}}
\newcommand{\filip}[1]{{\color{purple} \sf $\diamondsuit \diamondsuit \diamondsuit$ Filip: [#1]}}

\title{Classification of torsion of elliptic curves over quartic fields}
%\author{filip.najman }
\date{\today}


\begin{document}


\begin{abstract}
Let $E$ be an elliptic curve over a quartic field $K$. By the Mordell-Weil theorem, $E(K)$ is a finitely generated group. We determine all the possibilities for the torsion group $E(K)_\tor$ where $K$ ranges over all quartic fields $K$ and $E$ ranges over all elliptic curves over $K$. We show that there are no sporadic torsion groups, or in other words, that all torsion groups either do not appear or they appear for infinitely many non-isomorphic elliptic curves $E$.

Proving this requires showing that numerous modular curves $X_1(m,n)$ have no non-cuspidal degree $4$ points. We deal with almost all the curves using one of 3 methods: a method for the rank 0 cases requiring no computation (\Cref{prop_main}); the Hecke sieve, a local method requiring computer-assisted computations; and the global method, an argument for the positive rank cases also requiring no computation. We deal with the handful of remaining cases using ad hoc methods.



\end{abstract}

\author[Derickx]{Maarten Derickx}
\address{Maarten Derickx, University of Zagreb, Faculty of Science, Department of Mathematics, Bijeni\v{c}ka Cesta 30, 10000 Zagreb, Croatia}
\email{\url{maarten@mderickx.nl}}
\urladdr{\url{https://www.maartenderickx.nl/}}


\author[Najman]{Filip Najman}
\address{Filip Najman, University of Zagreb, Faculty of Science, Department of Mathematics, Bijeni\v{c}ka Cesta 30, 10000 Zagreb, Croatia}
\email{\url{fnajman@math.hr}}
\urladdr{\url{https://web.math.pmf.unizg.hr/~fnajman/}}
\thanks{The authors were supported by the  project “Implementation of cutting-edge research and its application as part of the Scientific Center of Excellence for Quantum and Complex Systems, and Representations of Lie Algebras“, PK.1.1.02, European Union, European Regional Development Fund and by the Croatian Science Foundation under the project no. IP-2022-10-5008.}
\subjclass[2020]{11G05, 14G05, 11G18}
\keywords{Modular curves, torsion, elliptic curves}



\maketitle







\section{Introduction}
Let $E$ be an elliptic curve over a number field $K$. The set $E(K)$ of $K$-rational points on $E$ has, by the Mordell-Weil theorem, the structure of a finitely generated abelian group. Hence it is isomorphic to $E(K)_\tor \times \Z^r$, where $E(K)_\tor$ is the \textit{torsion subgroup} of $E(K)$, and $r$ is the \textit{rank} of $E$ over $K$. The question of what are the possible values of $E(K)_\tor$ for $E$ running through all elliptic curves over all number fields $K$ of degree $d$ has received considerable attention. The first major result was Mazur's torsion theorem \cite{mazur77}, which, building on work of Kubert \cite{Kubert76} (and others), completed the classification of possible torsion groups of elliptic curves over $\Q.$ For quadratic fields, the classification was completed by Kamienny \cite{kamienny92}, building on work of Kenku and Momose \cite{KM88}. The last degree $d$ for which all the possible torsion groups of elliptic curves over degree $d$ number fields are known is $d=3$. The result was proved by Derickx, Etropolski, van Hoeij, Morrow and Zureick-Brown in \cite{Deg3Class}, building on previous work of Jeon, Kim and Schweizer \cite{JeonKimSchweizer04}, and Bruin and Najman \cite{BruinNajman15}.

The approach in all the mentioned papers is to study degree $d$ points on the modular curves $X_1(m,n)$. When $m > 4$ or $n > 4$ these are fine moduli spaces of (generalized) elliptic curves of triples $(E,P,Q)$ with $P$ and $Q$ being points on $E$ generating a subgroup isomorphic to $\tg{m}{n}$.

In this paper, we complete the classification of possible torsion groups of elliptic curves over quartic fields. Our main result is the following theorem.
\begin{theorem}\label{main_theorem}
 If $K$ varies over all quartic number fields and $E$ varies over
all elliptic curves over $K$, the groups that appear as
$E(K)_{\tor}$ are exactly the following
\begin{alignat*}{2}
& \torz{n},  && n=1-18, 20,21,22,24, \\
& \tg{2}{2n}, \quad \quad && n=1-9, \\
& \tg{3}{3n}, \quad && n=1-3, \\
& \tg{4}{4n}, \quad && n=1,2, \\
& \tg{5}{5}, \quad && \\
& \tg{6}{6}. \quad && \\
\end{alignat*}
\end{theorem}
One additional reason that the classification of possible torsion groups over quartic fields might be particularly interesting is that it seems it might be the last degree where there are no non-cuspidal \textit{sporadic points} on $X_1(m,n)$. Recall that a sporadic point on a curve $X$ is a point of degree $d$ such that there are only finitely many points of degree $\leq d$ on $X$. There are no non-cuspidal sporadic points on $X_1(m,n)$ of degree $1$ and $2$, while there exist sporadic cubic points on only one such curve, $X_1(21)$, and all of them correspond to a single elliptic curve (see \cite{najman16}). Jeon, Kim and Park \cite{JeonKimPark06} proved that all the groups in \Cref{main_theorem} appear infinitely often (see \Cref{infinite_quartic}), so it follows that there are no sporadic points of degree $4$ on $X_1(m,n)$. Moreover, from van Hoeij's results \cite{vanHoeij} one can easily see that there exist sporadic points of degree $5\leq d\leq 13$. See \Cref{sec:sporadic} for more details. Similarly, Clark, Genao, Pollack and Saia \cite{CGPS22} showed that $X_1(n)$ has sporadic points (corresponding to CM elliptic curves) for all $n>720$. These results show that there is an abundance of sporadic points on $X_1(n)$, and it might well be possible that $d=4$ is the largest degree such that, for all $n$, $X_1(n)$ has no degree $d$ sporadic points.

We use several tools to prove \Cref{main_theorem}. To deal with the cases where $\rk J_1(m,n)(\Q)=0,$ we first prove \Cref{prop_main} which allows us to eliminate all the cases where the hypothetical $E/K$ with $\tg{m}{n}\subseteq E(K)$ has bad reduction at all places above some rational prime $p$, e.g. when $Y_1(m,n)(\F_{p^d})= \emptyset$ for all $1\leq d\leq 4.$ To deal with the case where $E$ has good reduction at a prime over $p$, we develop the \textit{Hecke sieve} to show that all points in $Y_1(m,n)^{(4)}(\F_{p})$ cannot be reductions of points in $Y_1(m,n)^{(4)}(\Q)$ by studying the actions of Hecke operators. We give more details at the beginning of \Cref{sec:rank0}. 




To deal with the $\rk J_1(m,n)(\Q)>0$ cases we use the \textit{global method} (see \Cref{prop:hecke_gon}) and its variants to deal with almost all the remaining cases. The idea of this method is that, under some mild assumptions which are satisfied in all our cases, from a hypothetical point of degree $4$ on $Y_1(m,n)$ one could construct a map to $\PP^1$ of degree smaller than known gonality bounds, which is clearly impossible. 

A notable advantage of \Cref{prop_main} and the global method is their minimal reliance on computer calculations, especially in comparison to related works (e.g. \cite{Deg3Class}). These methods significantly reduce the number of cases requiring computational effort. For instance, we handle $X_1(45)$ using \Cref{prop_main} and $X_1(121)$ through the global method, whereas previous approaches demanded extensive computations to address these cases (see \cite{Deg3Class}).

%Additionally, our approach circumvents the use of formal immersions, which have been pivotal in classifying torsion groups over number fields of degrees $d = 1, 2, 3$. This difference suggests that our methods may offer greater potential for tackling problems over higher-degree number fields.



%A nice feature of \Cref{prop_main} and the global method is that they require almost no computer calculations, especially when compared to related papers in the field (e.g \cite{Deg3Class}). These two methods allow us to bring down the number of cases where we need to do computer calculations to a minimum. For example, we deal with the case of $X_1(45)$ using \Cref{prop_main}, and $X_1(121)$ using the global sieve, while prior work required extensive computation to deal with these cases (see \cite{Deg3Class}). Furthermore, our approach avoids the use of formal immersions, which have played a significant role in the classification of torsion groups over number fields of degrees $d = 1, 2, 3$. This distinction suggests that our methods may be better suited for addressing the problem in the context of higher-degree number fields.



Another important advantage of our work is that we do not use formal immersions, which have been used extensively in the classification of torsion groups over number fields of degree $d=1,2,3.$  The formal immersion criterion (without additional explicit computations) can only deal with elliptic curves that have bad reduction at all primes in the number field $K$ above some fixed rational prime $p$. So once the order $N$ of the point one wants to rule out is smaller than the Hasse bound at $2$ (in the best case), i.e. smaller than $(2^{d/2}+1)^2$, then the formal immersion approach runs into difficulties. On the other hand, we expect the lower bound for $N$ that the global method can deal with is quasi-linear in the level (more about this will be said in upcoming work by Derickx and Stoll). All this means that our methods are much more suitable to tackling the same question over higher degree number fields. As a demonstration, we apply our methods to higher degree number fields in \Cref{sec:higher_deg}. Specializing to quintic number fields, we eliminate more than half of the torsion groups one would need to eliminate to obtain a complete classification. We do this by "pure thought", i.e., by just checking the LMFDB and performing calculations that can easily be done by hand. 

In \Cref{sec:hecke} we work out a version of the Eichler-Shimura relation for $X_1(2,2n)$. This relation is already well-known for $X_1(n)$. Doing this was essential to be able to apply the Hecke sieve directly on $X_1(2,20)$ and $X_1(2,24)$ to deal with the non-cyclic cases $\tg{2}{20}$ and $\tg{2}{24}$.


To be able to use the Hecke sieve, we implemented the action of Hecke operators on non-cuspidal points on modular curves in \texttt{magma} \cite{magma}. 



%\todo[inline]{Say something about other methods if we end up using them.}

\subsection{Overview} \label{sec:overview}
For a finite abelian group $T$ we will say it is a \textit{quartic torsion group} if there exists an elliptic curve $E$ over a quartic field $K$ such that $E(K)_\tor\simeq T$. 

We now consider the groups that need to be eliminated to prove \Cref{main_theorem}. This is the list of all abelian groups $G$ with at most 2 generators such that:
\begin{itemize}
    \item $G$ occurs only finitely often as a quartic torsion group,
    \item for every proper subgroup $H \subsetneq G$, the group $H$ occurs infinitely often as a quartic torsion group,
    \item if $G$ contains a subgroup isomorphic to $(\Z/n\Z)^2$, then $\varphi(n)|4.$
    \item $|G|$ is divisible only by primes $\leq 17.$
\end{itemize}
The groups $G$ that occur infinitely often as quartic torsion groups can be found in \Cref{infinite_quartic}. The third bullet point is necessary because of the Weil pairing, while the third follows from \cite[Theorem 1.2.]{DKSS}.

The list of these groups is as follows:
\begin{itemize}
\item $\Z/n\Z, n =  {25^e, 26^e, 27^c, 28^e, 30^e, 32^e, 33^c, 34^e, 35^c, 36^e, 39^c, 40^a, 42^e, 44^a,}$  ${45^c, 48^a, 49^c, 51^c, 55^c, 63^f, 65^f, 77^d, 85^d, 91^d, 119^c, 121^d, 143^d, 169^d, 187^d,}$
${221^d, 289^d, }$

\item $\tg{2}{2n}$ $n=10^e,11^c,12^e,$
\item $\tg{3}{3n}$ $n=4^b,5^f,6^b,7^c,$
\item $\tg{4}{4n}$ $n=3^b,4^b,$
\item $\tg{5}{5n}$ $n=2^b,3^b,$
\item $\tg{8}{8n}$ $n=1^b.$
\end{itemize}

The cases have been marked according to the strategy they have been dealt with:
\begin{itemize}
    \item[$a$] These do not need to be studied separately, since there is $\Q$-rational map $X_1(4n) \to X_1(2,2n)$. So, the approach here is to show that $\tg{2}{2n}$ does not occur for $n=10,11$ and $12$ which is necessary in any case.
     
    \item[$b$] These were already done previously; see \Cref{thm:BN}.
    \item[$c$] These are done in \Cref{cor:easy} using \Cref{prop_main}.

    \item[$d$] These are done in \Cref{prop:global_argument} using a global argument.
    \item[$e$] These are done using the Hecke sieve in \Cref{cor:2_2n}, with an additional argument for $\torz{30}$ in \Cref{prop:30}.
    \item[$f$] These cases are solved using some ad hoc arguments in \Cref{prop:3x15,prop:63,prop:65}.
\end{itemize}

\begin{remark}
    While writing the paper, we became aware that some of the torsion groups $\torz{n}$ for which $\rk J_1(n)(\Q)=0$ have been eliminated, independently of us, by Michael Cerchia and Alexis Newton.  The values of $n$ they solved are $n=25, 26, 27, 28, 34, 35, 40$ (private communication).
\end{remark}




\subsection{Data availability and reproducibility}
All the code for the computations in this article is available at: 

\begin{center} \url{https://github.com/nt-lib/quartic-torsion} \end{center}

Claims in this article that are based on computations have been labeled in the source code in order to make it easy to navigate between this text and the code.
 For example, in the proof of \Cref{prop:30} there is a sentence: ``This shows that $z$
    is not the reduction of a point $x^{(4)}\in Y_1^{(4)}(30)(\Q),$ see \github{X1_30.m}{Claims 2,4,5} in the cases where $D$ is of degree $2$ and $3$ or irreducible of degree $4$."
 
 In order to verify, for example, Claim 2, one can open the file \githubbare{X1_30.m} and search for ``Claim 2" to find the following code that supports this claim.
{\small
 \begin{verbatim}
// Claim 2: Cannot have reduction of the form deg2 + 2 cusps
sieved := HeckeSieve(X, q, NoncuspidalPlacesUpToDiamond(X,2));
assert #sieved eq 0;
print "Claim 2 successfully verified";
\end{verbatim}
 }

 To be able to use the Hecke sieve, we implemented the action of Hecke operators on non-cuspidal points on modular curves over finite fields in \texttt{magma} \cite{magma}. The code for computing these Hecke operators and performing the Hecke sieve has been written in a reusable way and has been released as part of mdmagma \todo{release magma} at \begin{center} \url{https://github.com/koffie/mdmagma}. \end{center} \todo{release mdmagma v2}

The computations have been run on a server at the University of Zagreb with an Intel Xeon W-2133 CPU @ 3.60GHz with 12 cores and 64GB of RAM running Ubuntu 18.04.6 LTS and \texttt{Magma V2.28-3} \cite{magma}.
The computations took approximately 6.5 CPU hours and the maximal memory usage was roughly 1.5GB of RAM. Almost 3 of these 6.5 hours were spent on ruling out the possibility of $\tg{2}{24}$ torsion. The peak RAM usage was in ruling out $\tg{2}{20}$ torsion, with every other computation requiring less than half the amount of RAM.
More detailed logs of the computations are available in the \githubbare{logs} subdirectory of the quartic-torsion repository. \todo{fix this with ANTS template}

 \section*{Acknowledgements}
 We would like to thank Andrew Sutherland for computing the explicit equations for $X_1(m,n)$ that we used, and Barinder S. Banwait, Abbey Bourdon and Michael Stoll for helpful comments on an earlier version of the paper. 
 

\section{Conventions and notation}


Let $m\mid n$. We define $X_1(m,n)$ to be the moduli space parametrizing triples $(E,P,Q)$, where $E$ is a generalized elliptic curve, and $P$ and $Q$ are points generating a subgroup of $E[n]$ isomorphic to $\tg{m}{n}$. It is a smooth projective geometrically irreducible curve over $\Z[\zeta_m, \frac{1}{n}]$. It is a fine moduli space if $n>4$.

Cusps of the modular curve $X_1(m,n)$ correspond to \textit{N\'eron} $k$-gons, where $m\mid  k\mid n$ together with a level structure; for the definition see \cite[Sections 1.3 and 2.2.1]{DerickxPhD}. Over an algebraically closed field $K$ a N\'eron $k$-gon is a group scheme isomorphic to $\mathbb G_m\times \torz{k}$. If $n$ is invertible in $K$ then the level structure is an injective homomorphism $\phi:\tg{m}{n}\to \mathbb G_m(K)\times \torz{k}$ such that $\phi$ is surjective on the second coordinate.



\subsection{Analytic and algebraic ranks of modular abelian varieties}
If $A$ is an abelian variety over a number field $K$ then we use $\rk A(K)$ to denote the rank of the Mordell-Weil group $A(K)$. The integer $\rk A(K)$ is also called the algebraic rank of $A$. Similarly we use $\rk_{an} A(K)$ to denote the analytic rank of $A$ which is defined to be the order of vanishing of $L(A,s)$ at $s=1$. The BSD conjecture predicts that $\rk A(K) = \rk_{an} A(K)$. We use $S_k(\Gamma_1(n))^{new}$ to denote the set of all cuspidal newforms of weight $k$ for the congruence subgroup $\Gamma_1(n)$. If $f \in S_2(\Gamma_1(n))^{new}$ then we denote by $A_f$ the isogeny factor of $J_1(n)$ corresponding to $f$. This $A_f$ is a simple abelian variety over $\Q$ but might not be simple over $\overline \Q$. Note that for two newforms $f \in S_2(\Gamma_1(n))^{new}$ and $f' \in S_2(\Gamma_1(n'))^{new}$ we have $A_f=A_{f'}$ if and only if $n=n'$ and $f$ and $f'$ are Galois conjugates. It is known that every simple isogeny factor of $J_1(n)$ is isogenous to an abelian variety $A_f$ where $f \in S_2(\Gamma_1(n'))$ and $n' | n$. If $f$ is a newform in $f \in S_2(\Gamma_1(n))^{new}$ then we use $L(f, s)$ to denote its $L$-function and define $\rk _{an} f$ as the order of vanishing of $L(f,s)$ at $s=1$. If $f \in S_2(\Gamma_1(n))^{new}$ then we use $K_f \subset \C$ to denote the subfield generated by its coefficients. The field $K_f$ is known to be a number field isomorphic to $\Q \otimes \End A_f$. There is the following relation between $L$-functions: 
$$L(A_f,s) = \prod_{\sigma: K_f \to \C} L(\sigma(f),s).$$
Assuming the Birch-Swinnerton--Dyer conjecture, the order of vanishing of $L(\sigma(f),s)$ is the same for all choices of $\sigma \in \Hom_\Q (K_f, \C)$. In particular $\rk_{an} A_f(\Q) = [K_f: \Q] \rk_{an} f$.

A part of the BSD conjecture is known to be true for abelian varieties of the form $A_f$ due to the following celebrated result \cite[Corollary 14.3]{Kato2004} (we specialize by taking $\chi$ to be the trivial character):
\begin{theorem}
Let $A$ be an abelian variety over $\Q$ such that there is a surjective
homomorphism $J_1(n)\rightarrow A$ for some $n\geq 1$. Assume $L(A,1)\neq 0$. Then  $A(\Q)$ is finite.
\end{theorem}
This result of Kato generalizes earlier results by Kolyvagin and Logachev for abelian varieties attached to newforms in $S_2(\Gamma_0(n))$.

We use $\overline \Q$ to denote the algebraic closure of $\Q$ inside $\C$. Since the coefficients of newforms all lie in $\overline \Q$ we get an action of $\Gal(\overline \Q/\Q)$ on $S_2(\Gamma_1(n))^{new}$. We use $S_2(\Gamma_1(n))^{new}_{\Gal}:=  S_2(\Gamma_1(n))^{new}/\Gal(\overline \Q/\Q) $ to denote the set of all Galois orbits of newforms under this action. Since the isogeny class of $A_f$ only depends on the Galois conjugacy class of a newform, it also makes sense to define $A_{[f]}$  when $[f] \in S_2(\Gamma_1(n))^{new}_{\Gal}$ is a Galois orbit. In doing so, we get an isogeny decomposition 
$$J_1(n) \sim \bigoplus_{n' | n} \bigoplus_{f \in S_2(\Gamma_1(n'))^{new}_{\Gal}} A_{[f]}^{\sigma_0(n/n')},$$
where $\sigma_0(n/n')$ denotes the sum of positive divisors of $n/n'$.

\begin{example}
In this example we explain how to use the LMFDB to explicitly find a large rank 0 quotient of $J_1(N)$. The LMFDB contains an upper bound on the analytic rank of all weight 2 modular forms of level $\leq 1000$. In this article we often need the existence of a large isogeny factor $J_1(N)$ of algebraic rank 0 for $N<1000$. Such isogeny factors are easy to find using the data in the LMFDB. The \href{https://www.lmfdb.org/ModularForm/GL2/Q/holomorphic/?level_type=divides&level=65&weight=2&showcol=char_order.analytic_rank}{LMFDB data} on the modular forms of weight $2$ and level dividing $65$ shows that the Galois orbit of newforms with LMFDB label $\lmfdbnewform{65}{2}{a}{a}\in S_2(\Gamma_0(65))^{new}$ is the only orbit for which the analytic rank could potentially be positive.  In particular, it follows that every isogeny factor of $J_1(65)/J_0(65)$ has analytic and hence algebraic rank 0. Implying that $J_1(65)/J_0(65)(\Q)$ is a finite abelian torsion group. In the digital form of this article the words \href{https://www.lmfdb.org/ModularForm/GL2/Q/holomorphic/?level_type=divides&level=65&weight=2&showcol=char_order.analytic_rank}{LMFDB data} will always be a hyperlink that brings one to the LMFDB page containing the data that is needed for the argument. Similarly, LMFDB labels of newforms like \lmfdbnewform{65}{2}{a}{a} will also directly link to the LMFDB page on that newform.
\end{example}

\subsection{Hecke operators} \label{sec:hecke}
Let $q\nmid n$ be a prime, and $T_q$ be a Hecke operator on $X_1(n)$ defined as in \cite[Section 3.3]{DiamondIm}:
\begin{equation}\label{hecke_def_1}
    T_q(E,P)=\sum_{G \subseteq E[q] \text{ of order $q$}}(E/G, P \hspace{-0.25cm}\mod G).
\end{equation}
Recall that $T_{q,\F_q}$ acts on $J_1(n)_{\F_q}$ as $\Frob_q+\diamondop{q}_{\F_q,*}\Ver_q$, where Verschiebung $\Ver_q$ is the dual of the Frobenius, see \cite[(10.2.3)]{DiamondIm}. This is the Eichler-Shimura relation. 

For $q\nmid 2n$ a prime, we define the action of $T_q$ on $X_1(2,2n)$ similarly as in \eqref{hecke_def_1}:
\begin{equation}\label{hecke_def_2}
    T_q(E,P,Q)=\sum_{G \subseteq E[q] \text{ of order $q$}}(E/G, P \hspace{-0.25cm}\mod G, Q \hspace{-0.25cm}\mod G).
\end{equation}
We will need to consider $T_q$ as an operator on $X_1(2,2n)$. There is a degree 2 morphism 
$$X_1(4n)\rightarrow X_1(2,2n)=X_1(4n)/\diamondop{2n+1},$$
$$(E,P)\mapsto (E/\diamondop{2nP}, Q, P  \text{ mod}\diamondop{2nP}), $$
$\text{where }Q \text{ generates } E[2] /\diamondop{2nP},$ which is compatible with $T_q$, $\Frob_q$ and $\Ver_q$. It follows that $T_{q,\F_q}$ acts on $J_1(2,2n)_{\F_q}$ and also satisfies the Eichler-Shimura relation (as does $J_1(n)_{\F_q}$) 
\begin{equation}T_{q,\F_q}=\Frob_q+\diamondop{q}_{\F_q,*}\Ver_q, \label{eq:Eichler-Shimura}\end{equation} where $\diamondop{q}$ acts on $X_1(2,2n)$ as 
$$\diamondop{q}(E,P,Q)=(E,qP,qQ)=(E,P,qQ).$$

The following two propositions are not used in eliminating any of the groups in the paper, but we believe that they are interesting in their own right, and might be useful in future work. 

\begin{proposition} \label{prop:power-eichler-shimura}
Let $m=1$ or $2$ and $m|n$. Let $q\nmid n$ be a prime, $k>0$ an integer, and let $\Delta(n):=(\torz n)^\times/\diamondop{-1}$ denote the group of diamond operators. Then there exists a monic polynomial $\Psi_k(x) \in \Z[\Delta(n)][x]$ of degree $k$ such that $$\Psi_k(T_{q,\F_q})=\Frob_q^k+(\diamondop{q}_{\F_q,*}\Ver_q)^k.$$
%Here we see $\Z[(\torz n)^\times/\pm 1]$ as the group ring generated by the diamond operators $\diamondop{q}_{\F_q,*}$.
\end{proposition}
\begin{proof}
Let $a :=\Frob_q$ and $b := \diamondop{q}_{\F_q,*}\Ver_q$. Then the elementary symmetric polynomials on $a$ and $b$ are $a+b = T_q$ and $ab = q\diamondop{q}_{\F_q,*}$. Now 
\[T_q^k - a^k-b^k = ab\sum_{i=1}^{k-1} \binom k i a^{i-1}b^{k-i-1}.
\]
Now $\sum_{i=1}^{k-1} \binom k i a^{i-1}b^{k-i-1}$ is a symmetric polynomial of degree $k-2$ and can therefore be expressed in terms of $T_q$ and $q\diamondop{q}_{\F_q,*}$, where the power of $T_q$ occurring is at most $k-2$.  In particular, one can write $\sum_{i=1}^{k-1} \binom k i a^{i-1}b^{k-i-1} = \Psi'_k(T_q)$ for some polynomial $\Psi'_k \in \Z[\Delta(n)][x]$ of degree at most $k-2$. So $\Psi_k := x^k -q\diamondop{q}_{\F_q,*}\Psi'_k$ satisfies the proposition.
\end{proof}

\begin{example}
One can take $\Psi_1 =x$, $\Psi_2=x^2-2q\diamondop{q}$ and $\Psi_3=x^3-3q\diamondop{q}x$ in \Cref{prop:power-eichler-shimura}.
\end{example}

For a Hecke operator $T_q$ on $X_1(n)$ or $X_1(2,2n)$ we denote $A_q:=T_q-q\diamondop{q}-1.$
\begin{proposition} \label{prop:kill_tor}
Let $m=1$ or $2$ and $m|n$. If $q\nmid 2n$ is a prime that splits completely in a number field $K$ then
    \begin{equation}\label{eq:kill_tor}
    A_q(J_1(m,n)(K)_{\tor})=0.
\end{equation}
\end{proposition}
\begin{proof}
    Let $P$ be a point of order $k$ on $J_1(m,n)(K)$ and let $\wp$ be a prime of $K$ above $q$. Then $A_q(P)$ is a point of order dividing $k$. By the Eichler-Shimura relation \eqref{eq:Eichler-Shimura} and since $\Ver_q \circ \Frob_q=q$
    we have 
    $$T_{q,\F_q}(P_{\F_q})=\Frob_q(P_{\F_q})+\diamondop{q}_{\F_q,*}\Ver_q(q)=P_{\F_q}+q\diamondop{q} P_{\F_q}.$$
    The statement of the proposition now follows from the fact that reduction mod $\wp$ is injective on $J_1(m,n)(K)_{\tor}$ because $q>2$ and $\wp$ is unramified in $K$.
\end{proof}



\section{Previous results}
We recall previously proven results about torsion groups of elliptic curves over quartic fields. The following result tells us which quartic torsion groups appear infinitely often (in the sense that there exist infinitely many elliptic curves with different $j$-invariants with that torsion group over some quartic field). 

\begin{theorem}[{\cite[Theorem 3.6.]{JeonKimPark06}}]\label{infinite_quartic}
If $K$ varies over all quartic number fields and $E$ varies over
all elliptic curves over $K$, the groups which appear infinitely often as
$E(K)_{\tor}$ are exactly the following
\begin{alignat*}{2}
& \torz{n},  && n=1-18, 20,21,22,24 \\
& \tg{2}{2n}, \quad \quad && n=1-9 \\
& \tg{3}{3n}, \quad && n=1-3 \\
& \tg{4}{4n}, \quad && n=1,2 \\
& \tg{5}{5}, \quad && \\
& \tg{6}{6}. \quad && \\
\end{alignat*}
\end{theorem}




Some possible torsion groups have been ruled out in previous work of Bruin and Najman \cite{bruin_najman2016}, as can be seen from the following result.

\begin{theorem} \label{thm:BN}
  The following groups do not appear as $E(K)_{\tor}$ for any elliptic curve $E$ over any quartic field $K$:
\begin{alignat*}{2}
& \tg{3}{3n}, \quad && n\in \{4,6,9,11,13\}, \\
& \tg{4}{4n}, \quad && n\in \{3,4,7,11,13,17\}, \\
& \tg{5}{5n}, \quad && n\in \{2,3\}, \\
& \tg{8}{8}. \quad && \\
\end{alignat*}
\end{theorem}
\begin{proof}
    All the cases follow from {\cite[Theorem 8]{bruin_najman2016}}, except from $\tg{5}{5n}$ for $n=2$ and $3.$ These could a priori be torsion groups for elliptic curves only over $\Q(\zeta_5)$, but {\cite[Theorem 6]{bruin_najman2016}} shows that this is impossible.
\end{proof}




\section{Auxiliary results}
In this section we list (and slightly modify) some known results that we will use. 

\begin{theorem}[Abramovich's bound]
\label{abramovich}
    $$\gon_\C X_\Gamma>\frac{325}{2^{15}}[\PSL_2(\Z): \Gamma].$$
\end{theorem}
\begin{proof}
    This is \cite[Theorem 0.1]{abramovich} using the bound on the leading nontrivial eigenvalue of the non-Euclidean Laplacian $\lambda_1>\frac{975}{4096}$ from \cite[Appendix 2]{Kim03}.
\end{proof}

\begin{proposition}\label{prop:hecke-kernel}
Let $n$ be an integer, $q \nmid n$ a prime and let $H$ be a subgroup
  of~$(\Z/n\Z)^\times$ containing $-1$. Let $a=\sum_{i=1}^k a_i\diamondop{d_i} \in \Z[(\Z/n\Z)^\times/H]$ be a linear combination of diamond operators. We consider $t :=T_q - a$ 
as a correspondence on $X_1(N)$, inducing an endomorphism of the divisor group of $X_1(N)$
over $\C$. Then the kernel of $t$ is contained in the subgroup of divisors supported in cusps.
\end{proposition}
\begin{proof}
This is a slight generalization of \cite[Proposition 2.4]{DKSS} and the proof carries over ad verbatim.
\end{proof}


\begin{lemma}\label{lem:special_cusp}
Let $m=1$ or $2$, $n$ an integer such that $m|n$ and $q$ a prime that does not divide $2n$. Let $C_0 \in X_1(m,n)(\mathbb \Z[1/n])$ be the cusp whose moduli interpretation is the  N\'eron $n$-gon $E_0 := \mathbb G_m \times \torz{n}$ together with the points $P_0:=((-1)^{m-1},0)$ of order $m$ and $Q_0 :=(1,1)$ of order $n$. Then $A_q(C_0)$ is $0$.
\end{lemma}
\begin{proof}
    The proof uses the moduli interpretation of the Hecke operators. To this end, let $X_{0,1}(q,m,n)$ denote the modular curve parameterizing quadruples $(E,G,P,Q)$ where $E$ is a generalized elliptic curve, $G$ a subgroup of order $q$ and $P,Q$ points of order $m$ and $n$ such that $P,Q$ together generate a subgroup of order $mn$. Now consider the diagram
\[
\begin{tikzcd}
& X_{0,1}(q,m,n) \arrow[dl, "\pi_1"'] \arrow[dr, "\pi_2"] & \\
X_{1}(m,n) & & X_1(m,n)
\end{tikzcd}
\]
where $\pi_1(E,G,P,Q) = (E,P,Q)$ and $\pi_2(E,G,P,Q) =(E/G,P \mod G,Q \mod G)$. A small computation verifies that $T_q = \pi_{2,*} \circ \pi_1^*$ using the definition of $T_q$ as in \cref{hecke_def_1,hecke_def_2}. There are exactly two cusps $C_0'$ and $C_0''$ that map to $C_0$ under $\pi_1$. To be more precise, take $C_0'$ to be the cusp $(E_0,G_0',P_0,Q_0)$ with $E_0,P_0,Q_0$ as in the proposition and $G_0' = \mu_q \subseteq \mathbb G_m$, and take $C_0''$ to be the cusp $(E_0'',G_0'',P_0'',Q_0'')$ where $E_0'' := \mathbb G_m \times \torz{qn}$, $G_0'' = m\Z/qm\Z \subseteq \Z/qm\Z$, $P_0'':=((-1)^{m-1},0)$ and $Q_0'' :=(1,q)$. Note that $\pi_1$ is unramified at $C_0'$ and completely ramified of degree $q$ at $C_0''$ so that $\pi_1^*(C_0) = C_0'+qC_0''$. Now $\pi_2(C_0')=C_0$, since taking the quotient by $G_0'=\mu_q$ is the same as raising to the power $q$ on $\mathbb G_m$.  Similarly, one easily verifies $\pi_2(C_0'')=(E_0,P_0,pQ_0) = \diamondop{q}C_0$. Putting this all together yields $$T_q(C_0)=(\pi_{2,*} \circ \pi_1^*)(C_0)=C_0+q\diamondop{q} C_0$$ from which the proposition follows.

\end{proof}

The following proposition is a generalisation of a proposition of the first author \cite[Proposition 7.1]{DKSS}, which will appear in the unpublished article \cite{derickxstoll}. As that paper is still unpublished we include the proof here as well for completeness.


\begin{proposition}[(Derickx)] \label{prop:hecke_gon}
  Let $d \ge 1$, let $n$ be an integer, $q \nmid n$ a prime and let $H$ be a subgroup
  of~$(\Z/n\Z)^\times$ containing $-1$. We assume that there is an $a \in (\Z/n\Z)^\times/ H$ such that $T := (\diamondop{a}-1)(J_H(\Q))$ is finite.
  When $q=2$ we either assume that $\#A$ is odd or, more specifically, that the $2$-primary part of~$T$ is killed by $A_2$.
  We then set
  \[ k_q = \begin{cases}
           q+1 & \text{if $J_H(\Q)$ is finite, otherwise} \\
           2q + 1 & \text{if $a \in \set{q, q^{-1}}$,} \\
           2(q + 1)& \text{if $a \notin \set{q, q^{-1}}$.}
         \end{cases}
  \]
  Then $k_qd < \gon_{\Q}(X_H)$ implies that any point
  in~$X_H^{(d)}(\Q)$ without cusps in its support is a sum of orbits
  under~$\diamondop{a}$. In particular, the inequality above holds when
  \begin{equation}\label{DS-inequality}
      d < \frac{325}{2^{15}} \, \frac{[\SL_2(\mathbb Z) :\Gamma_1(n)]}{k_q \cdot \#H}.
  \end{equation} 
\end{proposition}
\begin{proof}
Let $c \in X_1(N)(\Q)$ be a rational cusp such that $A_q(c) = 0$; such a cusp exists by \Cref{lem:special_cusp}.  Let $c'$ be the image of $c$ in $X_H(\Q)$. Define 

\[ B_{a} :=\begin{cases}
           1 & \text{if $J_H(\Q)$ is finite, } \\
           (\diamondop{a}-1) & \text{otherwise.}
         \end{cases}
         \]
Let $D \in X_H^{(d)}(\Q)$ be an effective divisor of degree $d$ without cusps in its support, and consider the linear equivalence class $[D-dc']$ as a rational point in $J_H(N)(\Q)$. Then $A_qB_{a}([D-dc'])$ is $0$. Indeed, if $J_H(\Q)$ is finite then $A_qB_{a}([D-dc'])=0$ because $A_q$ kills all the torsion. If $J_H(\Q)$ is not finite, then $(\diamondop{a}-1)([D-dc'])$ is torsion by assumption and hence also killed by $A_q$.

Now since $A_qc' =0$ it follows that $A_qB_{a}(D-dc')=A_qB_{a}D$.
Now if $[A_qB_{a}D]=0$ there are two possibilities, either $A_qB_{a}D=0$ or there is a non-constant function $f$ such that $A_qB_{a}D = \ddiv f$.

First consider the case $A_qB_{a}D=0$. Since $D$ and hence $B_{a}D$ does not contain any cusps in its support, one can apply \Cref{prop:hecke-kernel} to deduce $B_{a}D=0$. If $J_H(\Q)$ is finite, this implies $D=0$, so in this case there does not exist an effective divisor $D \in X_H^{(d)}(\Q)$ of degree $d$ without cusps in its support, and there is nothing to prove. If $J_H(\Q)$ is not finite, then $B_{a}D=0$ implies $\diamondop a D =D$ forcing $D$ to be a sum of orbits under $\diamondop a$.

What remains to show is that the second case, $A_qB_{a}D = \ddiv f$ for some non-constant function $f$, cannot happen.

By separating the positive and the negative terms one can rewrite $A_qB_{a}D$ as \begin{align}A_qB_{a}D=\begin{cases}
           T_qD - (q\diamondop{q}+1)D  & \text{if $J_H(\Q)$ is finite, } \\
           (T_q\diamondop{a}+q\diamondop{q}+1)D - (T_q +\diamondop{a}(q\diamondop{q}+1))D & \text{otherwise}. 
         \end{cases}\label{eq:linear-equivalence}\end{align}
First consider the case when $J_H(\Q)$ is finite. Let $k_q:=q+1$. Note that $T_qD$  and $(q\diamondop{q}+1)D$ are both effective divisors of degree $k_qd$, so $\deg f \leq k_qd$ which contradicts $k_qd < \gon_\Q(X_H).$

Now consider the case when $J_H(\Q)$ is not finite. Let $k_q:=2(q+1)$.
If $a \notin \set{q,q^{-1}}$ then $(T_q\diamondop{a}+q\diamondop{q}+1)D$ and $(T_q +\diamondop{a}(q\diamondop{q}+1))D$ are effective divisors of degree $k_qdd$ and $\deg f \leq k_qd$ again contradicting $k_qd < \gon_\Q(X_H).$
In the cases $a=q$ and $a=\diamondop{q^{-1}}$ the right-hand side of \eqref{eq:linear-equivalence} simplifies to $(T_q\diamondop{q}+(q-1)\diamondop{q}+1)D - (T_q +q\diamondop{q^2})D$ respectively $ (T_q\diamondop{a}+q\diamondop{q})D - (T_q +(q-1)\diamondop{q}+\diamondop{q^{-1}})D.$ This simplification decreases the degree of the positive and negative part by 1 and hence there is again a contradiction with $k_qd < \gon_\Q(X_H)$ as before.
         
\end{proof}

\begin{corollary}
    Let the assumptions be as in \Cref{prop:hecke_gon}, and assume that the order of $a \in (\Z/n\Z)^\times/ H$ does not divide $d$. Then any rational point 
  on~$X_H^{(d)}$ is a sum of cusps and CM points.
\end{corollary}

\begin{theorem}[{Waterhouse, \cite[Theorem 4.1]{waterhouse69}}] \label{waterhouse}
Let $E$ be an elliptic curve over a finite field $\F_q$, where $q=p^a$. Let $\beta=q+1-\# E(\F_q)$. 
Then $\beta$ satisfies one of the following:
\begin{enumerate}
    \item $(\beta,p)=1$;
    \item $a$ is even and $\beta=\pm 2\sqrt q$;
    \item $a$ is even, $p\not\equiv 1 \pmod 3$ and $\beta=\pm \sqrt q$;
    \item $a$ is odd, $p=2$ or $3$ and $\beta=\pm p^{\frac{a+1}2}$;
    \item $a$ is odd and $\beta=0$;
    \item $a$ is even, $p\not\equiv 1 \pmod 4$ and $\beta=0$.
\end{enumerate}
\end{theorem}

\begin{corollary} \label{cor:curves_over_fp}

    \begin{enumerate}
        \item Suppose $mn$ is odd and $mn>25$. Then $X_1(m,n)$ has no non-cuspidal points over $\F_{2^d}$ for $d \leq 4.$

       \item Let 
        $$S_3=\left\{67,70,76,79,85,88,94,97\right\}.$$
        Suppose  $m \mid n$, $3\nmid mn$ and either $mn>100$, $mn\in S_3$, or $mn$ has no multiple $k$ such that $64\leq k\leq 100$. Then $X_1(m,n)$ has no non-cuspidal points over $\F_{3^d}$ for $d \leq 4.$
    \end{enumerate}
\end{corollary}
\begin{proof}
    This follows directly from \Cref{waterhouse}.
\end{proof}



\section{The Hecke sieve for the rank 0 cases} \label{sec:rank0}
Let $E$ be an elliptic curve over a number field $K$ of degree $d$ and $p \in \Z$ be a suitable prime. Our main approach is to split the problem up into cases depending on the reduction of $E$ above the different primes of $K$ above $p$. The argument exposition is tailored to the case in which $\rk J_1(m,n)(\Q(\zeta_m))=0$, but it could be generalized to $\rk J_1(m,n)(\Q(\zeta_m))>0$. We first prove \Cref{prop_main}, which is a modification of \cite[Theorem 1]{bruin_najman2016}, and which allows us to rule out the cases where $E$ has bad reduction at all primes above $p$; which is automatically satisfied if $Y_1(m,n)(\F_{p^d})=\emptyset$ for $1\leq d \leq 4.$  

To deal with the case in which $E$ has good reduction at least one prime of $K$ over $p$ we use a method that we call the \textit{Hecke sieve} (\Cref{prop:hecke_sieve,cor:bad_reduction}). The Hecke sieve works by studying the point $x^{(d)} \in X_1(m,n)^{(d)}(\Q(\zeta_m))$ corresponding to $\tg{m}{n} \hookrightarrow E(K)$ and using Hecke operators to rule out the different possibilities for $x^{(d)}_{\F_p}$ that have at least one non-cuspidal place in their support. The Hecke sieve has some similarities to the Mordell-Weil sieve, but has the advantage that one does not need to calculate any generators of the Mordell-Weil group $J_1(m,n)(\Q(\zeta_m))$.






\begin{proposition}\label{prop_main}
  Let $m\mid n$ be positive integers such that $mn>4$, $p$ be a prime not dividing $mn$, and $d$ an integer. Suppose $\rk J_1(m,n)(\Q(\zeta_m))=0$ and suppose one of the following is true:
  %Suppose that there exists no elliptic curve with a subgroup isomorphic to $\tg{m}{n}$ over $\F_{p^{fk}}$ for all $k\leq d$, where $f$ is the inertia degree of a prime $\wp'$ in $\Q(\zeta_m)$ above $p$, and that $\rk J_1(m,n)(\Q(\zeta_m))=0$. 
  %Suppose that for every cusp $c$ of $X_1(m,n)$ whose reduction is defined over $\F_{p^k}$ for some $k\leq d$, there is a single prime above $p$ in the field of definition $\Q(c)$ of $c$. 
  
\begin{itemize}
    \item[a)] If $p>2$ then $\gon_{\Q(\zeta_m)} X_1(m,n)>d$.

    \item[b)] If $p=2$ then either \begin{itemize}
        \item[i)]  $\gon_{\Q(\zeta_m)} X_1(m,n)>2d$ or
        \item [ii)] $\gon_{\Q(\zeta_m)} X_1(m,n)>d$, and there is a field $L \supseteq \Q(\zeta_m)$ and a prime $\wp''$ of $L$ above $p$ such that every cuspidal $\F_{\wp''}$-rational divisor on $X_1(m,n)$ lifts to a $L$-rational cuspidal divisor on $X_1(m,n)$ and reduction mod $\wp''$ is injective on $J_1(m,n)(L)[2]$. %\maarten{we need this complicated second condition for 27, since we only know that $J_1(27)(\Q)$ contains no $2$-torsion and not that $J_1(27)(\Q(\zeta_{27}))$ contains no $2$-torsion.} 
    \end{itemize}
\end{itemize}  
  Then every elliptic curve $E$ over a degree $d$ extension $K$ of $\Q(\zeta_m)$ with torsion $\tg{m}{n}\subseteq E(K)$ has good reduction at least one prime of $K$ above $p$.
  
  %  are no elliptic curves with torsion $\tg{m}{n}$ over a degree $d$ extension of $\Q(\zeta_m)$. 
  
  

\end{proposition}

\begin{proof}
    Let $\psi:\tg{m}{n} \hookrightarrow E(K)$ denote the inclusion. Fix a prime $\wp$ of $\Q(\zeta_n)$ above $p$ and above a prime $\wp'$ of $\Q(\zeta_m)$; we will denote by $\widetilde o$ the reduction of an object $o$ modulo $\wp$.  
    %Let $E$ be an elliptic curve over a number field $K$ of degree $d$ over $\Q(\zeta_m)$ with level structure $\psi:\tg{m}{n} \hookrightarrow E(K)$. 
    Let $x\in Y_1(m,n)(K)$ be a point representing $E$ together with the corresponding level structure $\psi$. Denote the corresponding points $x^{(d)}$ and $\tilde{x}^{(d)}$ on $X_1^{(d)}(m,n)(\Q(\zeta_m))$ and $ X_1^{(d)}(m,n)(\F_{\wp'})$ respectively.
    
     We first show that $E$ cannot have additive reduction at $\wp'.$ Assume for contradiction that $E$ has additive reduction modulo $\wp'.$ Then the identity component $E_{\F_{\wp'}}^0$ is of order $\#\F_{\wp'}$, a power of $p$, and $a:=[E_{\F_{\wp'}}:E_{\F_{\wp'}}^0]\leq 4$, so $\#E_{\F_{\wp'}}=a\cdot \#\F_{\wp'}$. The reduction mod $\wp'$ is injective on $\tg{m}{n}\subseteq E(K)[n]$, but $mn$ does not divide $a\cdot\#\F_{\wp'}$, so this is a contradiction. 

    If $E$ has good reduction modulo some prime of $K$ above $p$, there is nothing to prove. So $E$ has multiplicative reduction at all primes of $K$ above $p$. Then $\tilde{x}^{(d)}$ has to be a sum of reductions modulo $\wp$ of $\Q(\zeta_n)$-rational cusps, i.e., $\tilde{x}^{(d)}=\widetilde C$ for some effective cuspidal divisor $C \in X_1^{(d)}(m,n)(\Q(\zeta_n))$.

    Additionally let $c_0 \in X_1(m,n)(\Q(\zeta_m))$ be the cusp whose moduli interpretation is that of a N\'eron $n$-gon and $\tg{m}{n} \to \mathbb G_m(\Q(\zeta_m)) \times \Z/n\Z$ is given by $(a,b) \mapsto (\zeta_m^a,b)$.

    We have
    $$[x^{(d)}-C]=[x^{(d)}-dc_0]+[dc_0 -C].$$
    Since $[x^{(d)}-dc_0]\in J_1(n)(\Q(\zeta_m))=J_1(n)(\Q(\zeta_m))_{\tor}$ and since $dc_0 -C$ is a cuspidal divisor, by Manin-Drinfeld $[dc_0 -C]$ is torsion, so we conclude that $[x^{(d)}-C]\in J_1(n)(\Q(\zeta_n))_{\tor}$.
    
    Suppose first $p>2$. Then reduction mod p is injective on $J_1(n)(\Q(\zeta_m))_{tors}$ by \cite[Appendix]{katz81}, since $p$ is unramified in $\Q(\zeta_n)$, so it follows that $[x^{(d)}-C]=0$. Hence we have 
    $$\dim_{\Q(\zeta_m)} H^0(X_1(n)_{\Q(\zeta_m)}, \mathcal O(x^{(d)}))=\dim_{\Q(\zeta_n)} H^0(X_1(n)_{\Q(\zeta_n)}, \mathcal O(x^{(d)}))\geq 2.$$
    This inequality contradicts $\gon_{\Q(\zeta_m)} X_1(n)> d$, completing the proof if a) is satisfied.

    
    
    Suppose now that $p=2$ and $\gon_{\Q(\zeta_m)} X_1(n)> 2d$. By \cite[Prop 2.4]{Parent00} and the fact that $p$ is unramified in $\Q(\zeta_n)$, the kernel of the reduction mod $2$ of $J_1(n)(\Q(\zeta_n))$ is either trivial or a group of exponent $2$.
    
 Hence it follows that since $[\tilde{x}^{(d)} - \widetilde{C}]=0$, the divisor class $[x^{(d)}-C]\in J_1(n)(\Q(\zeta_n))$ has to be contained in $J_1(n)(\Q(\zeta_n))[2]$. Hence $2[x^{(d)}-C]=0$, so by the same argument as above, we arrive at a contradiction with $\gon_{\Q(\zeta_m)} X_1(n)> 2d$, completing the proof if b.i)  is satisfied.

    Suppose finally that $p=2$ and that the assumptions of b.ii) are satisfied. By assumption, we can choose $C$ to be $L$-rational, so $[x^{(d)}- C]$ is a $L$-rational divisor as well. Since $\tilde{x}^{(d)}=\widetilde C$ it follows that $[\tilde{x}^{(d)}-\widetilde C]=0$. By the injectivity of the reduction mod $\wp''$ on $J_1(m,n)(L)_\tor$ (which follows from \cite[Appendix]{katz81} and our assumptions), it follows that  $[x^{(d)}- C]=0$. By the same argument as above, we arrive at a contradiction with $\gon_{\Q(\zeta_m)} X_1(n)> d$, completing the proof if b.ii) is satisfied.    
\end{proof}

The following lemma will be useful when applying \Cref{prop_main}.
\begin{lemma}\label{lem:main}
Let $m|n$ be integers and $p$ a prime not dividing $n$, and let $E$ be an elliptic curve over a number field $K$ of degree $d$ such that $\tg{m}{n}\subseteq E(K)$. Suppose that for all $d'\leq d$ there does not exist an elliptic curve $E'/\F_{p^{d'}}$ with a subgroup of $E'(\F_{p^{d'}})$ isomorphic to $\tg{m}{n}$. Then $E$ does not have good reduction at any prime of $K$ above $p$.
\end{lemma}
\begin{proof}
    This follows immediately from the fact that reduction modulo a prime above $p$ is injective on all the torsion of order coprime to $p$.
\end{proof}




\begin{lemma}\label{lem:cusp_action}
Let $m=1$ or $2$ and $m|n$, $q$ a prime not dividing $2n$. Suppose $\rk J_1(m,n)(\Q)=0$. 
Then for any $x^{(d)}\in X_1(m,n)^{(d)}$ we have $[A_q(x^{(d)})]=0$.
\end{lemma}
\begin{proof}
    Let $C_0\in X_1(m,n)$ be the cusp such that $A_q(C_0)=0$; it exists by \Cref{lem:special_cusp}. Let $y=[x^{(d)}-dC_0]\in J_1(m,n)(\Q)$. By \Cref{prop:kill_tor} we have 
    $$0=[A_q(y)]=[A_q(x^{(d)})]-[A_q(C_0)]=[A_q(x^{(d)})].$$
\end{proof}


\begin{proposition}[Hecke Sieve] \label{prop:hecke_sieve} Let $m=1$ or $2$ and $m|n$, $p$ a prime not dividing $n$, $q$ a prime with $q\nmid 2pn$ and $1 \leq d'\leq d$ be integers. Assume $\rk J_1(m,n)(\Q)=0$. Suppose that 
\begin{itemize}


    \item[i)] $D\in Y_1(m,n)^{(d')}(\F_p)$ is a non-cuspidal effective divisor of degree $d'$ such that
    $$[A_{q}(D)]\neq 0,$$ and

    \item[ii)] $C \in X_1(m,n)^{(d-d')}(\F_p)$ is a cuspidal effective divisor of degree $d-d'$, such that $[A_{q}(C)]= 0$,
\end{itemize}
where $A_q=T_q-\diamondop{q}q-1$ as in \cref{eq:kill_tor}. Then a point $z \in Y_1(m,n)^{(d)}(\Q)$ cannot reduce to $D+C \in X_1(m,n)^{(d)}(\F_p)$.
\end{proposition}
\begin{proof} We have that 
$$[A_q(z)]_{\F_p}=[A_q(C)]+[A_q(D)]\neq 0,$$
while on the other hand 
$$[A_q(z)]=0$$
by \Cref{lem:cusp_action}. This is obviously a contradiction since $0\in J_1(m,n)(\Q)$ does not reduce modulo $p$ to a non-zero point in $J_1(m,n)(\F_p)$. 
\end{proof}
\begin{remark}
The only essential property of $A_q$ used in the proof of \Cref{prop:hecke_sieve} is that when $\rk J_1(m,n)(\Q)=0$ then $A_q(J_1(m,n)(\Q))=0$. In the case $\rk J_1(m,n)(\Q)>0$ one can still use the Hecke sieve by replacing $A_q$ with another Hecke operator $t$ such that $t(J_1(m,n)(\Q))=0$
\end{remark}

\begin{corollary} \label{cor:bad_reduction} Let $m=1$ or $2$ and $m|n$, and $p$ a prime not dividing $n$. Suppose $\rk J_1(m,n)(\Q)=0$. Suppose that
\begin{itemize}


    \item[i)] for all $1 \leq d'\leq d$ and all $D\in Y_1(m,n)^{(d')}(\F_p),$ there exists a prime $q_D\nmid 2pn$ such that  
    $$[A_{q_D}(D)]\neq 0,$$
    and

    \item[ii)] for all $d'\leq d$ such that $ Y_1(m,n)^{(d')}(\F_p)\neq \emptyset$ and all effective cuspidal divisors $C$ over $\F_p$ of degree $d-d'$, we have  $[A_{q_D}(C)]= 0$ for all $q_D$ used in step i) with $\deg D=d'$,
\end{itemize}
where $A_q=T_q-\diamondop{q}q-1$. Then every elliptic curve $E$ over a degree $d$ number field $K$ with torsion $\tg{m}{n}\subseteq E(K)$ has bad reduction at all primes of $K$ above $p$.
\end{corollary}
\begin{proof}
Let $x^{(d)}\in Y_1(m,n)^{(d)}(\Q)$ be the effective degree $d$ divisor corresponding to $E$ such that $\tg{m}{n}\subseteq E(K)$. Then $x^{(d)}_{\F_p}$ can be decomposed as $x^{(d)}_{\F_p}=C+D$ where $C$ is the sum of cusps and $D$ is a sum of non-cusps. If $E$ has good reduction at a prime above $p$, then $\deg D \geq  1$. However, by \Cref{prop:hecke_sieve} this cannot happen since there is no $z \in Y_1(m,n)^{(d)}(\Q)$ specializing to $x^{(d)}_{\F_p}=C+D$.
\end{proof}

\begin{lemma} \label{lem:split_completely_cusp}
     Let $m=1$ or $2$ and $m|n$, $q$ a prime not dividing $2n$. Suppose $\rk J_1(m,n)(\Q)=0$. Let $K$ be the field of definition of a cusp $C$ of $X_1(m,n)$. Assume that $q$ splits completely in $K$. Then 
     $$[A_q(C)]=0.$$
     \end{lemma}
\begin{proof}
    %The condition that $q$ splits completely in $K$ implies that all $K$-rational cusps reduce to $\F_q$-rational cusps. %It follows by our assumptions on $q$ that 
    %$$J_1(m,n)(K)_\tor \hookrightarrow J_1(m,n)(\F_q).$$
    %Now applying \Cref{prop:kill_tor} shows that $A_q(J_1(m,n)(\F_q))=0$. 
    By \Cref{prop:kill_tor} we have $A_q(J_1(m,n)(K)_\tor)=0$. Let $d=\deg C$ and let $C_0$ be the cusp defined as in \Cref{lem:special_cusp}. We have $[C-dC_0]\in J_1(m,n)(K)_\tor$ by Manin-Drinfeld, so $[A_q(C-dC_0)]=0$. As $[A_q(dC_0)]=0$ by \Cref{lem:special_cusp}, it follows that $[A_q(C)]=0.$
\end{proof}

We are ready to start eliminating possible quartic torsion groups.


\begin{corollary}\label{cor:easy}
    The groups $\torz{n}$ for $n\in\{27, 33,35,39 ,45,49,51,55,119\}$, $\tg{2}{22}$ and $\tg{3}{21}$ are not quartic torsion groups. 
\end{corollary}
\begin{proof}
    For all the torsion groups in the statement of the corollary, and the primes $p$ which we will choose for each group, it follows by \Cref{cor:curves_over_fp} that, for all $k\leq 4$, there are no elliptic curves containing that torsion group over $\F_{p^k}$.
    For all $\torz{n}$ we consider, it follows that $\rk J_1(n)(\Q)=0$ by \cite[Theorem 3.1]{Deg3Class}, while $\rk J_1(m,n)(\Q)=0$ follows from \cite[Theorem 4.1]{DerickxSutherland} for the groups $\tg{m}{n}$ with $m>1$ that we consider.
    
    For $\torz{n}$, with $n=33,35,39 ,45,49,51,55,119$ we apply \Cref{lem:main} and \Cref{prop_main} case b.i), using $p=2$; the necessary assumption $\gon_\Q X_1(n)>8$ follows from \cite[Theorem 3]{derickxVH}. 
    
    For $\torz{27}$ we apply \Cref{lem:main} and \Cref{prop_main}, case b.ii) with $L=\Q$ and $p=2$, using the fact that $J(\Q)$ has no $2$-torsion (see \cite[Table2]{Deg3Class}). To show that all $\F_2$-rational cuspidal divisors lift to $\Q$-rational cuspidal divisors, we note that all cusps of $X_1(27)$ are defined over subfields of $\Q(\zeta_{27})$, in which 2 is completely inert. 

    In the case $\tg{2}{22}$ we apply \Cref{lem:main} and \Cref{prop_main} case a) with $p=3$. We have that $\gon_\Q X_1(27)>4$ follows from \cite[Theorem 3]{derickxVH}.

    For $\tg{3}{21}$ we apply \Cref{lem:main} and \Cref{prop_main}, case b.i, with $p=2$, to show that there are no degree $2$ points over $\Q(\zeta_m)$ (and hence no quartic points over $\Q$) on $X_1(3,21)$. We see that $\rk J_1(3,21)(\Q(\zeta_3))=0$ from \cite[Theorem 4.1]{DerickxSutherland} and use \Cref{abramovich} to get
    $$\gon_{\Q(\zeta_m)} X_1(3,21)\geq\gon_{\C} X_1(3,21) \geq 6,$$
    completing this case. 
    
    
    %Furthermore the prime $2$ is completely inert in $\Q(\zeta_{27})$ so the splitting condition is satisfied. 

    %The prime $2$ is inert in $\Q(\zeta_3)$ and splits into primes of degree $8$ in $\Q(\zeta_{17})$ and $\Q(\zeta_{51})^+$.
\end{proof}

\begin{lemma}\label{lem:CC} Let $X:=X_1(m,n)$, with $m=1$ or $2$ and $m|n$, let $p$ and $q$ be different primes not dividing $2n$ and $k=1$ or $2$. Suppose $C\in X^{(k)}(\F_p)$ is a cuspidal divisor of degree $k$. Assume that 
\begin{itemize}
    \item[a)] if $k=1$, then $q$ splits completely in the field of definition of all cusps of $X$ that reduce to $\F_p$-rational cusps. 
    \item[b)] if $k=2$, then $q$ splits completely in the field of definition of all cusps of degree $1$ or $2$ on $X_{\F_p}$ which are reductions of cusps of higher degree on $X$.
\end{itemize}
    Then $[A_q(C)]=0.$
\end{lemma}
\begin{proof}
First consider the $k=1$ case. In this case $C$ is just a point, and is the reduction of a cusp $c$ (i.e., $c_{\F_p}=C$) defined over some number field $K$. If $q$ splits completely in $K$, we have $[A_q(c)]=0$ by \Cref{lem:split_completely_cusp} and hence $[A_q(C)]=0$. 

Now consider the $k=2$ case, and let $C\in X^{(2)}(\F_p)$ be a cuspidal divisor. If $C=C_1+C_2$ is reducible then $q$ by assumption splits in the fields of definition of the cusps of $X$ that reduce to $C_1$ and $C_2$, and $[A_q(C_1)]=[A_q(C_2)]=0$ by the same arguments as in the $k=1$ case. 

Let $c$ be an irreducible cuspidal divisor of $X$ such that $C$ is (a factor of) the reduction of $c$. If $\deg c>\deg C$, then by our assumptions $q$ splits in the field of definition of the points in $c$, and we conclude $[A_q(C)]=0$ by the same argumentation as in the $k=1$ case. Assume now $\deg c=\deg C$, (i.e. $c_{\F_p}=C$), and let $C_0$ be the cusp defined in \Cref{lem:special_cusp}. Then $[c-(\deg c) C_0]\in J(X)(\Q)_\tor$, so $[A_q(c)-(\deg c) C_0)]=0$. It follows that $[A_q(c)=0]$ and hence $[A_q(C)=0]$.
   
\end{proof}


We now apply the Hecke sieve to several torsion groups.

\begin{proposition}\label{cor:2_2n}
    The groups $\torz{n}$ for $n=25,26,28,32,24,36,42$, $\tg{2}{20}$ and $\tg{2}{24}$ are not quartic torsion groups.
\end{proposition}
\begin{proof}

We apply the Hecke Sieve to show that these groups are not quartic torsion groups. Let $T$ be one of the torsion groups from the statement of the proposition. For all the modular curves $X:=X_1(n)$ or $X_1(2,2n)$ (and defining similarly $Y:=Y_1(n)$ or $Y_1(2,2n)$) corresponding to our groups, we have $\gon_\Q X>4$ by \cite[Theorem 2.6 and Proposition 2.7]{JeonKimPark06} and $\rk J(X)(\Q)=0$ by \cite[Theorem 4.1]{DerickxSutherland} and \cite[Theorem 3.1]{Deg3Class}.

Let $p$ be a prime not dividing $\#T$, and hence a prime of good reduction of $X$. We can conclude by \Cref{prop_main} a) that for an elliptic curve $E$ over a quartic field $K$ with torsion $T \hookrightarrow E(K)$ has good reduction at a prime $\wp$ of $K$ over $p$.

So for the remainder of the proof we consider the case that $E$ has good reduction at least one prime $\wp$ of $K$ over $p$. We use $x \in X(K)$ denote the point corresponding to $E$ with the inclusion $T \hookrightarrow E(K)$, and let $x^{(4)} \in X(\Q)$ denote the effective divisor which is the sum of the Galois conjugates of $X$. The fact that $E$ has good reduction at $\wp$ is equivalent to $x^{(4)}_{\F_p}$ being of the form $D+C$ where $D \in Y^{(d')}(\F_p)$ is an effective divisor of degree $d'$ with $1 \leq d' \leq 4$ and $C$ a, possibly trivial, effective divisor supported on the cusps. 

We will complete the proof by verifying that for every $T$ we can find $p$ and $q$ such that the assumptions of \Cref{cor:bad_reduction} are satisfied. This will imply that $E$ has bad reduction modulo $\wp$, which will be a contradiction. In \Cref{table:Hecke_sieve} below we list the choices we make for our computations that verify this. 

In the 1st column we list the torsion groups $T$. In the second column we give the prime $p$ we use for this torsion group. In the columns $\deg d'$, for $d'=2,3,4$ we list the primes $q$ that we used to deal with that degree, if there is no value of $q$ listed this means $Y(\F_{p^{d'}})=\emptyset$. Note that  the column for $d'=1$ is omitted since for the $p$ we used we always have $Y(\F_{p})=\emptyset$. The irreducible divisors in $Y^{(d')}(\F_p)$ correspond to Galois orbits of points in $Y(\F_{p^{d'}})$ that do not come from a subfield. However $Y^{(d')}(\F_p)$ might also contain reducible divisors. Since $Y(\F_{p})=\emptyset$, the only way this can happen is if $d'=4$ and $D \in Y^{(4)}(\F_p)$ is the sum of two irreducible effective divisors of degree 2. This case is listed in the $\deg 2+2$ column. An absence of $q$ in this $\deg 2+2$ column denotes that $Y(\F_{p^2})=\emptyset$. If $Y(\F_{p^{d'}}) \neq \emptyset$ we explicitly check that $[A_q(D)]\neq 0$ for all irreducible effective divisors $x^{(d')}_{\F_p}\in Y^{(d')}(\F_p)$ in order to verify that \Cref{cor:bad_reduction} i) is satisfied. The value $q$ used  (or multiple values $q$ if more than one was used) is listed in the $\deg d'$ column. To explain the $\deg 2+2$ entry, note that we need to rule out the possibility that $x^{(4)}_{\F_p}\in Y^{(4)}(\F_p)$ is the sum of $2$ degree 2 effective irreducible divisors, so we check $[A_q(D)]\neq 0$ for all such divisors. For each $X_1(n)$ and $X_1(2,2n)$ one can find the computations we performed in our GitHub directory in the corresponding file \verb|X1_n.m| or \verb|X1_2_2n.m|.

It remains to verify \Cref{cor:bad_reduction} ii); this will in some cases impose congruence conditions on the choice of $q$ we are allowed to use. The only time we need to check anything is when $x^{(4)}_{\F_p}$ splits as $x^{(4)}_{\F_p}=D+C$, where $D$ is an irreducible effective non-cuspidal divisor and $C$ is an effective cuspidal (non-trivial) divisor. As $Y(\F_p)=\emptyset$ in all our cases, there are 2 cases we need to consider: when $\deg D=2$, and when $\deg D=3$. 

In the CC column in our table, a congruence condition on $q$ is listed, which guarantees that the assumptions of \Cref{lem:CC} are satisfied, and that we can conclude $[A_q(C)]=0$. If $\deg D=3$, this means that $q$ splits completely over the fields of definition of all cusps that reduce to cusps defined over $\F_p$. If $\deg D=2$, this means that $q$ splits completely in the field of definition of all cusps of degree $1$ or $2$ on $X_{\F_p}$ which are reductions of cusps of higher degree on $X$. We verify these claims in \github{congreuence_checks.m}{Claim 1} for $\deg D=3$ and in (see \github{congreuence_checks.m}{Claim 2}) for the $\deg D=2$ case.

%(see \github{congreuence_checks.m}{Claim 1}).

%First consider the $\deg D=3$ case. In this case $C$ is of degree 1, so just a point. If all $\F_p$-rational cusps are reductions of $\Q$-rational cusps, there is nothing to check. Indeed, let $c\in X(\Q)$ such that $c_{\F_p}=C$. By \Cref{lem:cusp_action} it follows that $[A_q(c)=0]$, and hence $[A_q(C)]=0$, proving our claim. Now consider the case where not all $\F_p$-rational cusps are reductions of $\Q$-rational cusps. However, $C$ will have to be the reduction of a cusp $c$ defined over some number field $K$ (where we can easily compute $K$).


%Now, if $q$ splits completely in $K$, we we have $[A_q(c)]=0$ by \Cref{lem:split_completely_cusp} and hence again $[A_q(C)]=0$.

%In the CC column in our table a congruence condition is listed which guarantees that $q$ splits completely over the fields of definition of all cusps that reduce to cusps defined over $\F_p$ (see \github{congreuence_checks.m}{Claim 1}). Our $q$ in the $\deg 3$ column will always be chosen to satisfy this condition. If this condition is empty, this means that all $\F_p$-rational cusps are reductions of $\Q$-rational cusps.

%\todo{modify this part of the proof.}

 %Now consider the $\deg D=2$ case. All of our $p$ are chosen so that the degree $1$ and $2$ cusps of $X_{\F_p}$ are reductions of degree $1$ and $2$ cusps of $X_\Q$ (see \github{congruence_checks.m}{Claim 2}). First suppose that all the cusps $C$ of $X_{\F_p}$ of degree $1$ and $2$ over $\Q$ reduce to cusps of the same degree over $\F_p$. Then we have $[A_q(C)]=0$ by the same argument as in the $\deg D=3$ case. In some cases it also happens that some degree $2$ cusps of $X_\Q$ reduce to $\F_p$-rational cusps of $X_{\F_p}$. In this case we choose $q$ which splits in the field of definition of such cusps $c$ of $X$. Then we have $[A_q(c)]=0$ by \Cref{lem:split_completely_cusp} and hence again $[A_q(c_{\F_p})]=0$.

%All the computations in this table can be found in 
%\todo{put the computations from here in a folder}.

\begin{table}[H]
\centering
\begin{tabular}{|c|c|c|c|c|c|c|}
\hline
$(m,n)$ & $p$ & \multicolumn{5}{|c|}{$q$} \\
\hline
& &  CC & $\deg2$ &   $\deg 3$ & $\deg 2+2$ & $\deg 4$  \\
\hline
$(1,25)$ & $3$ &    &    &   &   & $7$ \\
\hline
$(1,26)$ & $7$ & & $3$ &  $3$ &  $3$     & $3$ \\
\hline
$(1,28)$ & $5$ & $ 1 \pmod 4$ & $13$  & $13$ &  $3$ &   $3$ \\
\hline
$(1,32)$ & $3$ & &  &    $5$  &  & $5$ \\
\hline
$(1,34)$ & $3$ &  &   &   &    & $5$ \\
\hline
$(1,36)$ & $5$ &  $ 1 \pmod 4$ &   &$13$ & & $7$ \\
\hline
$(1,42)$ & $11$ & & $5$ & $5$ & $5$ &     $5$ \\
\hline
$(2,20)$ & $7$ & & $3$  & $3$  &  $3$    & $3$ \\
\hline
$(2,24)$ & $5$ & $ 1 \pmod 4$ &   &     $17$ & &  $7,11$ \\
\hline


\end{tabular}
\caption{}\label{table:Hecke_sieve}
\end{table}
\end{proof}

\begin{remark}
    In the proof of \Cref{cor:2_2n} we used explicit equations for $X_1(2,2n)$. We used optimized models obtained using the methods of \cite[\S 3]{sutherlandtorsion} to make the run time finish in a reasonable amount of time.
\end{remark}


The case of the torsion group $\torz{30}$ requires some additional argumentation. 

\begin{proposition} \label{prop:30}
    The group $\torz{30}$ is not a quartic torsion group. 
\end{proposition}
\begin{proof}We largely follow the proof of \Cref{cor:2_2n}.  We have $\gon_\Q X_1(30)$ $>4$ by \cite[Theorem 2.6]{JeonKimPark06} and $\rk J_1(30)(\Q)=0$ by \cite[Theorem 3.1]{Deg3Class} so that we can apply \Cref{prop_main} a) to show that any elliptic curve over a quartic field $K$ with $\torz{30}$ torsion has to have good reduction at least one prime $\wp$ of $K$ over $p$.

    We apply the Hecke sieve for $X_1(30)$ using $p=7$ in order to rule out the possible reductions $Y_1^{(4)}(30)(\F_7)$, this time using \Cref{prop:hecke_sieve} directly instead of \Cref{cor:bad_reduction}. Let  $z\in Y_1^{(4)}(30)(\F_7)$, and write $z=D+C$, where $D$ is a non-trivial effective non-cuspidal divisor and $C$ is an effective cuspidal divisor. As $Y_1(30)(\F_7)=\emptyset$, $D$ is of degree $2,3$ or $4$.

    We apply the Hecke sieve using $q=13, 13$ and $11$, respectively, for degrees $2, 3$ and $4$, respectively. This shows that $z$
    is not the reduction of a point $x^{(4)}\in Y_1^{(4)}(30)(\Q),$ see \github{X1_30.m}{Claims 2,4,5} in the cases where $D$ is of degree $2$ and $3$ or irreducible of degree $4$, respectively.
    

    However, there are 2 points $z',z''\in Y_1^{(4)}(30)(\F_7)$ (up to the action of diamond operators), each of which is the sum of 2 irreducible degree $2$ non-cuspidal divisors, and which are not eliminated by the Hecke sieve. 

    To eliminate them, we note that $J_1(30)(\Q)$ is generated by differences of cusps \cite[Corollary 4.14]{Deg3Class}. Let $C\in X_1(30)(\Q)$ be some rational cusp. It follows that if either $z'$ or $z''$ were reductions modulo $7$ of $x^{(4)}\in Y_1^{(4)}(30)(\Q)$, then $[z'-4C_{\F_7}]$ or $[z''-4C_{\F_7}]$ would have to lie in the subgroup of $J_1(30)(\F_7)$ generated by sums of cusps that lift to $\Q$. However, using an explicit computation we check that $[z'-4C_{\F_7}]$ or $[z''-4C_{\F_7}]$ are not in the subgroup generated by cusps \github{X1_30.m}{Claim 3}, completing the proof.
\end{proof}


\begin{comment}
    We will study the reduction modulo the prime $p=7$ for the torsion groups $\tg{2}{20}$ and the prime $p=5$ for $\tg{2}{24}$.

     Since $\gon_\Q X_1(2,20)>4$ and $\gon_\Q X_1(2,24)>4$ by \cite[Proposition 2.7]{JeonKimPark06} and $\rk J_1(2,20)(\Q)=\rk J_1(2,24)(\Q)=0$ by \cite[Theorem 4.1]{DerickxSutherland}, we can conclude by \Cref{prop_main} a) that for an elliptic curve $E$ over a quartic fields $K$ with torsion $\tg{2}{20}\subseteq E(K)$ ($\tg{2}{24}\subseteq E(K)$, respectively) have good reduction at a prime $\wp$ over $7$ ($5$, respectively).


    We now check that the assumptions of \Cref{prop:bad_reduction} are satisfied, which will imply that $E$ has bad reduction modulo $\wp$, which will be a contradiction.  

    We have that $Y_1(2,20)(\F_{7^{d'}})$ has $0,16,128,704$ Galois orbits of points for $d'=1,2,3,4,$ respectively \github{X1_2_20.m}{Claim 1}. The curve $Y_1(2,24)(\F_{5^{d'}})$ has $0,0,16, 96$ Galois orbits of points for $d'=1,2,3,4,$ respectively \github{X1_2_24.m}{Claim 1}. We explicitly check that $[A_q(D)]\neq 0$ for all effective divisors corresponding to these orbits (see \github{X1_2_20.m}{Claims 3--6} and \github{X1_2_24.m}{Claims 2--4}, respectively) which verifies that \Cref{prop:bad_reduction} i) is satisfied. For $X_1(2,20)$ we always use $q=3$, while in the case $X_1(2,24)$ we use $q=19$ for $d'=3$ and $q=7$ and $11$ for $d'=4$. 
    
    The cusps of $X_1(2,20)$ are all defined over subfields of $\Q(\zeta_{20})$. %The prime $7$ is inert in $\Q(\zeta_{20})^+$ as $7$ is of order $4$ in $\torz{20}^\times /\diamondop{-1}$, which 
    We explicitly check that for any cusp $C$ of $X_1(2,20)$ we have $[\Q(C):\Q]=[\F_7(C_{\F_7}):\F_7]$ (see \github{X1_2_20.m}{Claim 2}), and in particular this implies that all $\F_7$-rational cuspidal divisors $C_{\F_7}$ lift to $\Q$-rational cuspidal divisors. As $[A_q(C_\Q)]=0$ by \Cref{lem:cusp_action} (for all $q$), it follows that $[A_q(C_{\F_7})]=0$.

     In the case of $X_1(2,24)$ \todo{We need to take care of the cusps here.}, we have $Y_1(2,24)(\F_{5^{d'}})\neq \emptyset$ for only $d'=3,4$, so we need only show that $[A_q(C_{\F_5})]=0$ for the $\F_5$-rational cusps. The cusps of $X_1(2,24)$ are defined over subfields $\Q(\zeta_{24})^+=\Q(\sqrt 2, \sqrt 3).$ Since $5$ does not spilt completely in any subfield of $\Q(\sqrt 2, \sqrt 3)$ except $\Q(\sqrt 6)$, the only cusps that will reduce to $\F_5$-rational cusps are those that are defined over $\Q$ (there are $8$ of them) and those that are defined over $\Q(\sqrt 6)$ (there are $4$ of them).
     In the case $d'=4$ there is nothing to prove, so let $d'=3$. In this case we use $q=19$, and since $19$ splits in $\Q(\sqrt 6)$, the result now follows from \Cref{lem:split_completely_cusp}. 
    
    In the case $q=19$ for $d'=3$, which is the only computationally costly case, we note that all the $16$ orbits are in one $\Aut X_1(2,24)$-orbit, so it is enough to perform the computation to verify \Cref{prop:bad_reduction} i) for only one orbit.  
\end{comment}
    





\begin{comment}

\subsection{Direct analysis over \texorpdfstring{$\F_p$}{Fp}} \todo{check if we still want to include it given the work of DZB's students}
In this section we follow the strategy of \cite[Section 5.1]{Deg3Class}; following loc. cit. we call this method \textit{direct analysis over $\F_p$}. Let $p \nmid n$ be an odd prime, $X:=X_1(n)$, $J:=J_1(n)$, let $\red_X$ and $\red_J$ be the reductions of $X^{(4)}$ and $J$ modulo $p$. Assume $\gon_\Q X>4$ and $\rk J(\Q)=0$. We have the commutative diagram.
\begin{equation} \label{direct_diagram}
    \begin{tikzcd}[sep = large]
    X^{(4)}(\Q) \arrow[hook]{r}{\iota}  \arrow[hook]{d}{\red_X} &  J(\Q)  \arrow[hook]{d}{\red_J} 
    \\ 
    X^{(4)}(\mathbb{F}_p) \arrow{r}{\iota_p} & J(\mathbb{F}_p)
\end{tikzcd}\end{equation}
The map $\red_J$ is injective by \cite[Appendix]{katz81}; it follows that $\red_X$ must be injective. Note that $\iota_p$ will also be injective if $\gon_{\F_p} X>4$.

The set $X^{(4)}(\mathbb{F}_p)$ can be explicitly computed (in reasonable time). Hence the idea is to explicitly compute $\iota^{-1}(\red_J^{-1}(\img \iota_p))\subseteq X^{(4)}(\Q)$. \todo{include Maarten's speedup.}

It is also possible to use the prime $p=2$, however then $\red_J$ might not be injective. So, to apply this method, one needs to explicitly check whether $\red_J$ is injective on $J(\Q)$ for $p=2$.

\begin{proposition}
    The groups $\torz n$ for $n \in \{25,26,28\}$ are not quartic torsion groups.
\end{proposition}
\begin{proof}
    By direct analysis over $\F_3$ for $n=25,26, 28$. The code for this can be found in the Github repository \todo{cite}.
\end{proof}



\begin{proposition}
    $\torz{42}$ is not a quartic torsion group.
\end{proposition}
\begin{proof}\todo{check this proof, also do we still want to include it?}
    All points on $X_1(42)^{(4)}(\Q)$ reduce to points on $X_1(42)^{(4)}(\Q)$ supported on cusps. This follows from Maarten's computation showing that a Hecke operator does not send the divisor to 0 in the Jacobian.

    To show that the map is a formal immersion we need to show that it has gonality $\geq 5$ over $\F_5$.
    
    Let $X:=X_1(42)$. Suppose $\gon_{\F_5}X=d\leq 4$.    
    Since $X$ has a degree 3 map $f$ to the genus $5$ curve $X_1(21)$, the Castelnuovo-Severi inequality gives
    $$25=g(X)\leq 3\cdot 5 +2\cdot (d-1)\leq 21,$$
    which is a contradiction unless the map, call it $g$, to $\PP^1$ factors through $f$. But then it follows that $d\leq 2\deg f\leq 6$, since $X_1(21)$ is not hyperelliptic over $\F_5$ \cite[Section 2.1.]{DerickxNajman24}.
    
\end{proof}
\end{comment}
\subsection{The group $\tg{3}{15}$}

We show that $\tg{3}{15}$ is not a quartic torsion group by considering it over finite fields of different characteristics. 

\begin{proposition}\label{prop:3x15}
    The group $\tg{3}{15}$ is not a quartic torsion group.
\end{proposition}
\begin{proof}
    Let $X:=X_1(3,15)$; we consider this curve over $K:=\Q(\sqrt {-3})$. The curve $X$ has 64 cusps in total, all of which are defined either over $K$ or over $\Q(\zeta_{15})$.
    
     Let $C_K$ be the group generated by differences of $K$-rational cusps, and $C_{\F_4}$ the subgroup generated by differences of $\F_4$-rational cusps. We have $\#C_{\F_4}\leq \#C_{K}$, since all $\F_4$-rational cusps are reductions of $K$-rational cusps. We compute $J(\F_7)[2]\simeq (\torz{2})^3$ \github{X1_3_15.m}{Claim 1} and $C(\F_4)[2]\simeq (\Z/2\Z)^3$ \github{X1_3_15.m}{Claim 2}. Since 
    $$(\torz{2})^3\simeq C_{\F_4}[2] \twoheadleftarrow C_K[2] \hookrightarrow J(K)[2] \hookrightarrow J(\F_7)[2] \simeq (\torz{2})^3, $$
    we conclude that the reduction modulo the prime above $2$ is injective.  

    Now the case follows \Cref{prop_main} b.ii) and \Cref{lem:main} using $K=\Q(\zeta_3)$ and $d=2$ and the fact that all the points in $X(\F_4)$ and $X(\F_{16})$ are cusps.   
\end{proof}


\begin{remark} Let $C$ be the cuspidal subgroup of $J:=J_1(m,n)$. Note that $C(\Qbar)\leq J(\Qbar)_{\tor}$, by Manin-Drinfeld, and that any prime of good reduction of $J$ is unramified in $\Q(\zeta_n)$, in which all the cusps of $X_1(m,n)$ are contained. This implies that reduction $C(\Qbar)\rightarrow C(\overline \F_p)$ is injective for all primes $p>2$ not dividing $n$.

Let $C$ be the cuspidal subgroup of $J_1(3,15)$. We compute that 
$$C(\Q(\zeta_{15}))[2]\simeq(\torz{2})^6, \text{ while }C(\F_{16})[2]\simeq(\torz{2})^5,$$ which implies that the reduction modulo a prime above $2$ is not injective on $C(\Qbar)$. This is the first instance that we are aware of, where the reduction of the cuspidal subgroup modulo a prime of good reduction is not injective on a modular Jacobian. 
\end{remark}







\section{The positive rank cases}




\subsection{The global method}


We can apply \Cref{prop:hecke_gon} to eliminate most of the positive rank cases. 

\begin{proposition} \label{prop:global_argument}
    $\torz{n}$ is not a quartic torsion group for 
    $$n\in \{77, 85,91,121, 143,169, 187, 221, 289\}.$$
\end{proposition}
\begin{proof}
    For $n\in\{77,85,91, 121, 169, 221, 289\}$ we have $$\rk J_1(n)(\Q)=\rk J_0(n)(\Q),$$ so $A:=(\diamondop{a}-1)J_1(n)(\Q)$ is finite for any non-trivial element $a\in (\torz{n})^\times$. For $n\in\{143,187\}$ we have $\rk J_1(n)(\Q)>\rk J_0(n)(\Q)$. In both cases, the positive rank parts of $J_1(n)/J_0(n)$ correspond to a unique Dirichlet character Galois orbit. These character orbits are \lmfdcharorbit{143}{h} for $n=143$ (corresponding to the newform \lmfdbnewform{143}{2}{h}{a}) and \lmfdcharorbit{187}{g} for $n=187$ (corresponding to the newform \lmfdbnewform{187}{2}{g}{a}), and in both cases the characters in these orbits are all of order 5. For $n=143$ (respectively $187$) let $\chi$ be a character in the character orbits \lmfdcharorbit{143}{h} (respectively \lmfdcharorbit{187}{g}). If we choose $a$ in the kernel of $\chi$ then $A:=(\diamondop{a}-1)J_1(n)(\Q)$ is finite, and we can apply \Cref{prop:hecke_gon} as well. Choosing $a$ with $\diamondop{a}$ of order coprime to $5$ suffices to ensure $\chi(a)=1$. 
    
    
    We will use $q=3$ in all cases, and will have $k_3=7$ in all cases, apart from $n=121,143$ and $187$, where $k_3=8$. In \Cref{table:DS} we list the values $a \in (\torz{n})^\times $ that we use, the order of $\diamondop{a}$ and the value of  $b(n):=\frac {325} {2^{15}}\frac{[\SL_2(\mathbb Z) :\Gamma_1(n)]}{k_3 \cdot \#H}$; we need to have $b(n)> 4$ to proceed. The value $d_{CM}X_1(n) $ is the least degree of a CM point on $X_1(n)$, as listed in the GitHub repository \url{https://github.com/fsaia/least-cm-degree} of the paper \cite{CGPS22}.
    
\begin{table}[H]
\centering
\begin{tabular}{c|c|c|c|c}
$n$ & $a$ & $|\diamondop{a}|$ & $\lceil b(n)\rceil $ &  $d_{CM}X_1(n) $ \\ 
\hline
$77$  & $3$   & $30$ & $5$  & $60$  \\
$85$  & $3$   & $16$  & $5$  & $32$  \\
$91$  & $3$  & $6$  & $6$  & $24$  \\
$121$ & $56$  & $11$ & $10$ & $110$ \\
$143$ & $67$  & $12$ & $13$ & $120$ \\
$169$ & $3$  & $39$  & $21$ & $52$  \\
$187$ & $122$ & $16$ & $22$ & $160$ \\
$221$ & $3$  & $48$  & $35$ & $96$  \\
$289$ & $3$ & $272$  & $59$ & $136$ \\
\end{tabular}
\caption{}\label{table:DS}
\end{table}

We now apply \Cref{prop:hecke_gon}.  Since $b(n) > 4$, \eqref{DS-inequality} is satisfied. Thus all rational points in $D\in X_1^{(4)}(n)(\Q)$, not supported on the cusps, are invariant under the action of $\diamondop{a}$. If $D$ has no fixed points of $\diamondop{a}$ in its support, then it has to be of degree at least $|\diamondop{a}|$ (the order of $\diamondop{a}$), which is not possible by \Cref{table:DS}. This implies that the support of $D$ consists of points corresponding to elliptic curves with $j=0$ or $j=1728$. But from the aforementioned results \cite{CGPS22}, it follows that $X_1(n)$ has no CM points of degree $< d_{CM}X_1(n) $, which shows that this is impossible.



\end{proof}

\begin{remark} \label{rem:6.2}
    Note that in \Cref{prop:global_argument} we also prove results for higher degree number fields. What follows from our proof is that there are no elliptic curves over number fields of degree $n$ over number fields of degree $d\leq \lceil b(n)\rceil$  such that $|\diamondop{a}|$ does not divide $d$, for the values of $n$ in \Cref{table:DS}.
\end{remark}


\subsection{The groups $\torz{63}$ and $\torz{65}$} \label{sec:63_65}

We now deal with the cases $\torz{63}$ and $\torz{65}$ which turn out to be harder and require some ad hoc argumentation. %, the main reason being that $\rk J_1(n)(\Q)>0$ for $n=63,65$. 


\begin{lemma} \label{lem:cusps}Below we determine the cusps of $X_1(n),$ for $n=63, 65$, that are defined over $\F_{2^k}$  and $\F_{3^k}$ for $k\leq 4$.
    \begin{enumerate}
        \item $n=63$: The cusps corresponding to $63$-gons are defined over $\F_2$, the cusps corresponding to $21$-gons are defined over $\F_4$, the cusps corresponding to $9$-gons are defined over $\F_8$, while the remaining cusps are defined over $\F_{2^k}$ for some $k\geq 5.$
        \item $n=65$ in characteristic $2$:  The cusps corresponding to $65$-gons are defined over $\F_2$, and the cusps corresponding to $13$-gons are defined over $\F_{16}$, while the remaining cusps are defined over $\F_{2^k}$ for some $k\geq 5.$
        \item $n=65$ in characteristic $3$:  The cusps corresponding to $65$-gons are defined over $\F_3$, and the cusps corresponding to $13$-gons are defined over $\F_{81}$, the cusps corresponding to $5$-gons are defined over $\F_{27}$ while the $1$-gons are defined over $\F_{3^k}$ for some $k\geq 5.$
    \end{enumerate}
\end{lemma}
\begin{proof}
 The cusps corresponding to $n$-gons are defined over $\Z[1/n]$ and the $1$-gons over $\Z[1/n, \zeta_{n}+\zeta_{n}^{-1}]$. For the squarefree cases $n=pq$, the $p$-gons are defined over $\Z[1/n, \zeta_p]$. 
\end{proof}


\begin{proposition}\label{prop:63}
    $\torz{63}$ is not a quartic torsion group.
\end{proposition}

\begin{proof}
    We will eliminate this case by showing that there are no non-cuspidal quartic point on the genus $49$ curve $Y:=X_1(63)/\diamondop{8}$ which are images of quartic points of $X_1(63)$, by following the strategy of \Cref{prop_main}. The group of diamond operators on $Y$ is $(\torz{63})^\times/ \diamondop{\pm 8}\simeq (\torz{3})^2$. This implies that isogeny factors of $J(Y)$ correspond to weight 2 newforms of level dividing 63 and with character order $1$ or $3$. The \href{https://www.lmfdb.org/ModularForm/GL2/Q/holomorphic/?level_type=divides&level=63&weight=2&char_order=1%2C3&showcol=analytic_rank.char_order.prim&hidecol=analytic_conductor.field.cm.traces.qexp} {LMFDB data} shows that all these newforms have analytic rank 0 and hence $J(Y)(\Q)$ has algebraic rank $0$.

    
    
    

    Suppose $E/K$ is an elliptic curve over a quartic field $K$ with a point of order $63$ over $K$ and let $x_0=(E,P)\in X_1(63)(K)$ for some $P\in E(K)$ of order 63. Let $\wp$ be a prime of $K$ over $2$, and let $\wp'$ be a prime of $\overline K$ above $\wp$. Since there are no elliptic curves with a point of order $63$ over $\F_{2^k}$ for $k\leq 4$, it follows that $x_0$ reduces to a cusp $\widetilde c\in X_1(63)(O_K/\wp)$. There is a unique cusp $c_0\in X_1(63)(\overline K)$ that reduces modulo $\wp'$ to $\widetilde c_0$. %Since $\widetilde{x}$ is defined over $\F_{2^k}$ for some $k\leq 4$, it follows that $c_0$ cannot correspond to a $1$-gon, since their reductions are defined over $\F_{2^6}$.

  
    Let $f:X_1(63)\rightarrow Y$ be the quotient map. Let $c=f(c_0)$; this is a cusp as $f$ sends cusps to cusps. Let $x=f(x_0)\in Y(K)$. The points $x$ and $c$ reduce modulo $\wp'$ to the same point in $\overline{\F}_2$; moreover, they are both defined over some $\F_{2^k}$ for $k\leq 4$. %The cusps of $Y$ (excluding $1$-gons that we don not need to consider) are defined over $\Q$ ($63$-gons and $21$-gons), $\Q(\zeta_9)^+$ (7-gons), $\Q(\zeta_7)^+$ (9-gons). The prime $2$ is inert in all these fields.

    Let $D=\sum_{i=1}^4x^{\sigma_i}$ be the degree 4 divisor which is the sum of all the conjugates of $x$; it is a $\Q$-rational divisor. As in the proof of \Cref{prop_main} there is a rational cuspidal divisor $C$ such that it reduces to the same divisor modulo 2 as $D.$ 

    By \cite[Prop 2.4]{Parent00}, the kernel of the reduction mod $2$ of $J(Y)(\Q)$ is either trivial or a group of exponent $2$. Hence it follows that since $[\widetilde D - \widetilde{C}]=0$, the divisor class $[D-C]\in J(Y)(\Q)$ has to be contained in $J(Y)(\Q)[2]$. Hence $2[D-C]=0$, so there exists a degree $8$ morphism $Y\rightarrow \PP^1$. But \Cref{abramovich} gives $\gon_\Q Y \geq \gon_\C Y>8$, so we have arrived at a contradiction. 

\end{proof}









%This method can be applied when $\gon_\Q X> 4n$ and $\rk J_1(\Q)/J_0(\Q)=0$, where $n$ can always be taken to be $8$, and can be taken to be $6$ if the $t_2$ \todo{define this} operator kills all the $2$-torsion of $J_1$, and can be taken to be $4$ if $\rk J_1(\Q)=0$.

%\subsection{The group \texorpdfstring{$\Z/65\Z$}{Z/65Z}}

\begin{proposition}\label{prop:65}
    The group $\torz{65}$ is not a quartic torsion group.
\end{proposition}

\begin{proof}
We start similarly as in \Cref{prop_main}.  Fix a prime $\wp$ of $\Q(\zeta_{65})$ above $2$; we will denote by $\widetilde o$ the reduction of an object modulo $\wp$ or mod $2$ depending on whether the object is defined over $\Q(\zeta_{65})$ or over $\Q$.  Suppose for contradiction that $E$ is an elliptic curve over a number field $K$ of degree $4$ over $\Q(\zeta_{65})$ and such that $\torz{65} \subseteq E(K) $. Let $x\in Y_1(65)(K)$ be a point representing $E$ together with a point of order $65$. Consider the corresponding points $x^{(4)}$ and $\tilde{x}^{(4)}$ on $X_1^{(4)}(65)(\Q)$ and $ X_1^{(4)}(65)(\F_2)$ respectively. Note that the Hasse-Weil bound implies that for $1\leq i \leq 4$ there are no elliptic curves over $\F_{2^i}$ with a point of order $65$. This means that $\tilde{x}^{(4)}$ has to be a sum of reductions modulo $\wp$ of $\Q(\zeta_{65})$-rational cusps. In particular, we can write $\tilde{x}^{(4)}=\widetilde C$ for some effective cuspidal divisor $C \in X_1^{(4)}(65)(\Q(\zeta_{65}))$. 
    
By looking up the analytic ranks of all weight 2 newforms of level dividing 65, it follows that $J_1(65)(\Q)$ has rank 1. In fact the \href{https://www.lmfdb.org/ModularForm/GL2/Q/holomorphic/?level_type=divides&level=65&weight=2&showcol=char_order.analytic_rank}{LMFDB data} shows more precisely that:
$$\rk(J_1(65)(\Q))=\rk(J_0(65)(\Q))= \rk(J_0^+(65)(\Q))=1.$$ 

In particular, it follows that $(\diamondop{a}-1)J_1(65)(\Q) \subseteq J_1(65)(\Q)_{tors}$ for all $a \in (\torz{65})^\times$.

    Let $c_0 \in X_1(65)(\Q)$ be the cusp whose moduli interpretation is that of a N\'eron $65$-gon and and the point of order $65$ is given by $(1,1) \in \mathbb G_m(\Q) \times \Z/65\Z$. 

    We have
    $$(\diamondop{3}-1)[x^{(4)}-C]=(\diamondop{3}-1)[x^{(4)}-4c_0]+(\diamondop{3}-1)[4c_0 -C].$$
    Since $(\diamondop{3}-1)[x^{(4)}-4c_0]\in (\diamondop{3}-1)J_1(65)(\Q) \subseteq J_1(65)(\Q)_{\tor}$ and since $(\diamondop{3}-1)(4c_0 -C)$ is a cuspidal divisor, by Manin-Drinfeld it follows that $(\diamondop{3}-1)[4c_0 -C]=[(\diamondop{3}-1)(4c_0 -C)]$ is torsion, so we conclude that $(\diamondop{3}-1)[x^{(4)}-C]\in J_1(65)(\Q(\zeta_{65}))_{\tor}$.

    By \cite[Prop 2.4]{Parent00} and the fact that $2$ is unramified in $\Q(\zeta_{65})$, the kernel of the reduction mod $2$ of $J_1(n)(\Q(\zeta_{65}))$ is either trivial or a group of exponent $2$.

     Since $(\diamondop{3}-1)[\tilde{x}^{(4)} - \widetilde{C}]=0$, the divisor class $(\diamondop{3}-1)[x^{(4)}-C]\in J_1(65)(\Q(\zeta_{65}))$ has to be contained in $J_1(65)(\Q(\zeta_{65}))[2]$. Hence $2(\diamondop{3}-1)[x^{(4)}-C]=0$; we show that this cannot happen by following the proof strategy of \Cref{prop:hecke_gon}.
     
     There are two ways one could have $2(\diamondop{3}-1)[x^{(4)}-C]=0$. Either 
     \begin{enumerate}
         \item[(i)]  $2(\diamondop{3}-1)(x^{(4)}-C)=\ddiv(f)$ for some non-constant function $f \in \C(X_1(65))$, or
         \item[(ii)] $2(\diamondop{3}-1)(x^{(4)}-C)=0$ as a divisor.
    \end{enumerate}
    
    Case (i) cannot happen. Indeed, in this case the degree of $f$ is at most $\deg(2x^{(4)}+2\diamondop{3}C) = 16$, which contradicts the gonality bound $\gon_\Q X_1(65)\geq \gon_\C X_1(65)\geq 20$ obtained using Abramovich's bound. 
    
    Now we show that the case (ii) is impossible as well. The divisor $2(\diamondop{3}-1)x^{(4)}$ is supported at non-cuspidal points, while $2(\diamondop{3}-1)C$ is supported at the cusps. The equality $2(\diamondop{3}-1)(x^{(4)}-C)=0$ forces both $2(\diamondop{3}-1)x^{(4)}=0$ and $2(\diamondop{3}-1)C=0$. In particular, one gets $\diamondop{3}x^{(4)}=x^{(4)}$. This equality can happen only if $x$ is a fixed point of $\diamondop{3^4}$. Since $3^4$ is of order $3$ in $(\torz{65})^\times/\set{\pm 1}$ this means that $E$ has to have CM by $\Q(\zeta_3)$. This is contradicting the fact that the least degree of a CM point on $X_1(65)$ is 24, as listed in the GitHub repository \url{https://github.com/fsaia/least-cm-degree} of the paper \cite{CGPS22}.
\end{proof}

\begin{remark}
    The proof of \Cref{prop:65} combines the strategies of \Cref{prop:hecke_gon,prop_main}. It should be possible to extend \Cref{prop_main} with this extra case. However, the statement of \Cref{prop_main} is already rather verbose. So for readability reasons, we chose not to extend this proposition even further with another case.
\end{remark}

\section{Higher degree number fields} \label{sec:higher_deg}

In this section we briefly explore how our methods can be applied to ruling out possible torsion groups over high degree number fields. We only use ``pure thought", using only data from LMFDB and calculations that can be done by hand. 

\begin{theorem}\label{thm:d5r1}
Let $n$ be one of the values in Table \ref{table:higher_degree}. Then there does not exist an elliptic curve $E$ with a point of order $n$ over a number field $K$ of degree $d\leq \lceil b(n) \rceil$ such that $|\diamondop{a}|$ does not divide $ d$ .
\begin{table}[H]
\centering
\begin{tabular}{c|c|c|c|c|c}
$n$ & $|\chi|$& $a$ & $|\diamondop{a}|$ & $\lceil b(n)\rceil $ &  $d_{CM}X_1(n) $ \\ 
\hline
$95$ &$12$ & $3^{12}$  & $3$  & $6$  & $72$  \\
$119$ &$\emptyset$ & $3$  & $48$  & $10$  & $96$  \\
$125$ &$\emptyset$ & $3$  & $50$  & $10$  & $50$  \\
$133$ &$3,6$ & $3^6$  & $3$  & $11$  & $36$  \\
 $209$ &$5$ & $3^5$  & $18$  & $27$  & $180$  \\
 $247$ &$12,12$ & $3^{12}$  & $3$  & $38$  & $72$  \\
 $323$ &$8$ & $3^8$  & $18$  & $65$  & $288$  \\
 $361$ &$\emptyset$ & $3$  & $171$  & $93$  & $114$  \\
\end{tabular}
\caption{}\label{table:higher_degree}
\end{table}
\end{theorem}
The values $n$, $a$, $|\diamondop{a}|$, $\lceil b(n)\rceil $ and $d_{CM}X_1(n)$ are the same as in \Cref{table:DS}, while the values in $|\chi|$ represent the orders of characters corresponding to the factors of $J_1(n)/J_0(n)$ of positive rank. 

\begin{proof}
    The proof  follows along the same lines as the proof of \Cref{prop:global_argument}. We choose $a:=3^k$ where $k$ is a multiple of all values in $|\chi|$ to ensure that $a$ is in the kernel of $\chi$. Using the same arguments as in \Cref{prop:global_argument}, we conclude that if $|\diamondop{a}|$ does not divide $d$, $d<d_{CM}X_1(n)$ and $d\geq \lceil b(n)\rceil $, then there does not exist an elliptic curve $E$ with a point of order $n$ over a number field $K$ of degree $d$.
    \end{proof}

We now consider the cyclic torsion groups that need to be ruled out over quintic fields. We apply \Cref{prop_main} to eliminate some torsion groups whose corresponding modular curve has a rank 0 Jacobian over $\Q$.  

\begin{proposition}\label{prop:d5r0}
    There are no elliptic curves with a point of order $49,51,55$ or $75$ over degree $5$ number fields. 
\end{proposition}
\begin{proof}
    Let $n:=49,51,55$ or $75$. In all cases, we have $\rk J_1(n)=0.$ We apply \Cref{lem:main} and \Cref{prop_main} case b.i), using $p=2$; the necessary assumption $\gon_\Q X_1(n)>10$ follows from Abramovich's bound. In particular, we obtain $\gon_\Q X_1(n)\geq \lceil b(n) \rceil$ where $\lceil b(n) \rceil$ is as in the following table. 
    \begin{table}[H]
\centering
\begin{tabular}{c|c||c|c||c|c||c|c}
$n$ & $\lceil b(n)\rceil$& $n$ & $\lceil b(n)\rceil$&$n$ & $\lceil b(n)\rceil$&$n$ & $\lceil b(n)\rceil$ \\ 
\hline
$49$ & $12$ & $51$ & $12$ & $55$ & $15$ & $75$ & $24$ \\
\hline

\end{tabular}
\caption{}\label{table:gonalities}
\end{table}
    
\end{proof}

By \cite[Theorem 1.1.]{DerickxSutherland} the groups that appear infinitely often as quintic torsion groups are $\tg{m}{n}$ for 
\begin{equation}
    \{(1,n): 1\leq n\leq 25,\ n \neq 23 \} \cup \{(2,2n): 1\leq n\leq 8\}.
\end{equation}
To completely determine the possible torsion groups of elliptic curves over quintic fields, using the same reasoning as in \Cref{sec:overview} we would need to rule out the following 45 groups: 
\begin{itemize}


\item $\Z/n\Z, \quad n =  26, 27, 28, 30, 32, 33, 34, 35, 36, 38, 39, 40, 42, 44, 45, 48,$\\ $ 49, 50, 51,
55, 57, 63, 65, 75, 77, 85, 91, 95, 119, 121, 125, 133, 143, 169,$\\ $ 187,209, 221, 247, 289, 323, 361.$
\item $\tg{2}{2n}$, for $9\leq n\leq 12.$
\end{itemize}
 In \Cref{thm:d5r1} we eliminate $8$ groups, in \Cref{prop:d5r0} another $4$ groups, and in \Cref{rem:6.2} we have already eliminated $7$ groups. Note that the groups $\torz{4n}$ can be ruled out once the groups $\tg{2}{2n}$ are eliminated, due to the existence of a $\Q$-rational map $X_1(4n) \to X_1(2,2n)$. This allows us to ``eliminate" the groups $\torz{4n}$ for $9\leq n\leq 12.$ In total we have eliminated 23 out of 45 possible quintic torsion groups. 
 
 Finally, note that the cases $\torz{28}$ and $\torz{30}$ \textit{do} appear as quintic torsion groups (finitely many times), see \cite{vanHoeij} and \Cref{sec:sporadic}. So these torsion groups obviously do not need to be eliminated; one instead needs to find all the degree 5 non-cuspidal points on $X_1(28)$ and $X_1(30)$.





\begin{comment}
\section{The modular curves $X_1(2,2n)$}
It turns out that all the previous methods do not work for the curves $X_1(2,2n)$ for $n=10,12.$


\begin{proof}
    The curve $X_0(65)$ is of genus $5$, and 
    $$\rk(J_1(65)(\Q))=\rk(J_0(65)(\Q))= \rk(J_0^+(65)(\Q))=1.$$ \todo{justify this rank computation}Suppose there exists an elliptic curve $E$ over a quartic field $K$, together with a point $P\in E(K)$ of order $65$. Let $x\in X_1(65)(K)$ be the corresponding point. and let $D\in X_1^{(4)}(65)(\Q)$ be the trace of the point. Let $\pi:X_1(65)\rightarrow X_0(65)$ be the projection sending $(E,P)$ to $(E,\diamondop{P})$. Let $y:=\pi(x)$ and let $D_0 \in X_0^{(4)}(65)(\Q)$ be the trace of $y$. By \Cref{cor:curves_over_fp} and \Cref{lem:cusps}, $D_0$ reduces modulo $2$ to $4\cdot c_{\F_2}$ for some cusp $c$ of $X_0(65)$.
    
    We compute $J_0(65)(\F_2)\simeq \tg{2}{168}$ and $J_0(65)(\F_3)\simeq \tg{6}{84}$, and that the differences of cusps of $X_0(65)$ generate a subgroup isomorphic to $\tg{2}{84}$. From this we conclude that $J_0(\Q)_\tor\simeq \tg{2}{84}$, and that it is generated by differences of cusps. We compute that the reductions of cusps modulo 2 generate a subgroup of $J_0(65)(\F_2)$ isomorphic to $\torz{84}$, which implies that the reduction modulo 2 on $J_0(\Q)_\tor$ has kernel isomorphic to $\torz{2}$. 
 
    Since $(w_{65}-\Id)(D_{0,\F_2}) = (w_{65}-\Id)(4\cdot c_{\F_2})$, we have that $(w_{65}-\Id)(D_0-4c)$ is in the kernel of $\red_2$, and we conclude that 
    \begin{equation} \label{eq:65}
        2(w_{65}-\Id)(D_{0})=2(w_{65}-\Id)(4c).
    \end{equation}
    
    We have the following commutative diagram, with the vertical sequence being exact:
    \begin{equation} \label{65_diagram}
    \begin{tikzcd}[sep = large]
    & & \torz{2} \arrow{d}
    \\
    X_1(65)^{(4)}(\Q) \arrow{r}{\pi}   &  X_0(65)^{(4)}(\Q)  \arrow{d}{\red_2} \arrow{r}{w_{65}-\Id} &  J_0(65)^{(4)}(\Q)_\tor \arrow{d}{\red_2}
    \\ 
    & X_0(65)^{(4)}(\mathbb{F}_2) \arrow{r}{w_{65}-\Id} & J(\mathbb{F}_2)   
\end{tikzcd}\end{equation}

By \Cref{cor:curves_over_fp} we know that there exists an elliptic curve over $\F_{81}$ with a point of order $65$, hence there are non-cuspidal points on $X_1(65)(\F_{81})$. By a brute force computation, we check that all such points lie in a single orbit under diamond operators and Frobenius. We compute that all such curves have $j$-invariant $\pm i$, where $i^2=-1$. Suppose $D_{0,\F_3}$ is the trace of such a point. Then since reduction modulo 3 is injective on $J_0(65)(\Q)$, it follows that the \eqref{eq:65} should hold over $\F_3$. We check that explicitly; it does not hold, and this completes the proof. 
\todo{write about cusps over $\F_3$.}


\todo{explain that we have a formal immersion at the relevant cusps over $F_3$ and add what is necessary.}

    
\end{proof}
\end{comment}


\appendix
\begin{comment}
\section{Reduction of torsion}

This section appendix contains a generalization of \cite[Lemme 4.14]{Parent99} for which we couldn't find a reference.
\begin{lemma}Let $R$ be a discrete valuation ring whose fraction field $K$ has characteristic 0, and whose residue field $k$ has characteristic $p>0$. Let $\pi$ generate the maximal ideal of $R$. Let $G$ be a group scheme over $R$, $t$ a positive integer and $P \in G(R)$ a torsion point of order $n$ such that $P_{R/\pi^t R} = 0$. Then
\begin{itemize}
    \item There is an integer $k$ such that $n=p^k$
    \item 
\end{itemize}
\end{lemma}
\begin{proof}
    As in \cite[Lemme 4.14] let $\spec A$ be an affine open neighbourhood of $0_{k}$ of $G$. Note that $P$ induces a morphism of group schemes $ \torz{n}_R \to G$, that factorises via $\spec A$. Let $F_K$ be the image of $\torz{n}_K$ in $\spec A_K$, and $F$ it's schematic closure in $\spec A$. In  \cite[Lemme 4.14] it is proved that $F = \spec A/I$ is an affine groups scheme such that $A/I$ is a free $R$ module of rank $n$.

    Translation by $P$ induces an automorphism of $F$ and hence one of $A/I$. Since $A/I$ is a free $R$-module of rank $n$, this automorphism can be represented by a matrix in $\GL_n(R)$. Because $P_{R/\pi^t R} = 0$ this matrix is of the form $1 + \pi^t A$ with $A \in M_n(R)$.
\end{proof}

\begin{proposition}
    Let $M \in \GL_n(R)$ be a matrix of multiplicative order $n$ that reduces to the identity mod $\pi^t$. Then
    \begin{enumerate}
    \item There is an integer $k$ such that $n=p^k$
    \item 
    \end{enumerate}
    
\end{proposition}

\begin{proof}
For the proof part (i). Let $q$ be a prime dividing $n$. Then $M^{n/q}$ is of order $q$ and we can write $M^{n/q} = 1+\pi^{t'}A$ for some positive integer $t'$ and a matrix $A \in M_n(R)$ that has at least entry that is a unit.
Now $(1+\pi^{t'}A)^q=1$ so that 
\begin{align*}0 = (1+\pi^{t'}A)^q -1 \equiv q\pi^{t'}A \mod \pi^{t'+1}
\end{align*}
and $q\pi^{e'}A \equiv 0 \mod \pi^{t'+1}$ if and only if $p=q$. This proofs $(i)$.

For the second part write $M = 1+\pi^{t'}A$ for some positive integer $t' \geq t$ and a matrix $A \in M_n(R)$ that has at least entry that is a unit. By part $1$ we have $(1+\pi^{t'}A)^{p^k}=1$, it follows that
\begin{align*}0 = (1+\pi^{t'}A)^{p^k} -1 &= \binom{p^k} {1}\pi^{t'}A^1 + \binom{p^k} {2}\pi^{2t'}A^2 + \ldots \pi^{t'p^k}A^{p^k} 
\end{align*}
\end{proof}
Since
\end{comment}

\section{Sporadic points on \texorpdfstring{$X_1(n)$}{X1(n)} and \texorpdfstring{$X_0(n)$}{X0(n)}}
\subsection{Sporadic points on \texorpdfstring{$X_1(n)$}{X1(n)}}
\label{sec:sporadic}
In the introduction we mentioned that sporadic points of degree $5\leq d\leq 13$ have been found on $X_1(n)$ in \cite{vanHoeij}. However, loc. cit. only lists points of low degree, and it is not immediately clear from the data that these points are sporadic. We give a table from which it should be clear that our claim is true, and for which $n$ there are $X_1(n)$ with degree $d$ sporadic points. The data in the table about $\rk J_1(n)$ is obtained from \cite[Theorem 3.1]{Deg3Class}, the data about $\gon_\Q X_1(n)$ from \cite{derickxVH} and the low degree points themselves are collected from \cite{vanHoeij}. Recall that for a point $x$ of degree $n$ on a curve $X$, it is sufficient that either $n<\frac 1 2\gon_\Q X$ (see \cite{frey}), or $n<\gon_\Q X$ and $\rk J(X)(\Q)=0$ (see \cite[Proposition 2.3]{DerickxSutherland}).


\begin{table}[H]
\centering
\begin{tabular}{c|c|c|c}
$d$ & $n$ & $\rk J_1(n)$ & $\gon_\Q X_1(n)$  \\ 
\hline
$5$  & $28$   & $0$ & $6$    \\
$6$  & $37$   & $>0$  & $18$   \\
$7$  & $33$  & $0$  & $10$  \\
$8$ & $33$  & $0$ & $10$  \\
$9$ & $31$  & $0$ & $12$ \\
$10$ & $29$  & $0$  & $11$  \\
$11$ & $35$ & $0$ & $12$  \\
$12$ & $39$  & $0$  & $14$  \\
$13$ & $39$ & $0$  & $14$  \\

\end{tabular}
\caption{}\label{table:sporadic}
\end{table}

There are a number of points of degree $d=14,15, \ldots$, which are likely to be sporadic; however, lower bounds on the gonalities of $X_1(n)$ for $n\geq 41$ that are good enough to prove this have not been determined, and this is the obstacle of going any further than $d=13$ at the moment. 

\subsection{Sporadic points on \texorpdfstring{$X_0(n)$}{X0(n)}}
The evidence for the abundance of sporadic points on $X_0(n)$ of all degrees is even stronger than for $X_1(n)$, which leads us to make the following conjecture. 
\begin{conjecture}
For every positive integer $d$ there exists an $n$ such that there exists a sporadic point of degree $d$ on $X_0(n)$.\end{conjecture}

We now give some evidence for the conjecture. 


\begin{proposition} \label{sporadic}
    If $d\leq 2166$ or there exists an imaginary quadratic number field with class number $d$, then there exists a degree $d$ sporadic point on some $X_0(n)$. 
\end{proposition}

\begin{proof}
    Let $d$ be an integer, $K=\Q(\sqrt{-n})$, where $n$ is square-free, an imaginary quadratic field whose absolute value of the discriminant is $\Delta$ and which has class number $h_\Delta=d$. Let $E_\Delta$ be an elliptic curve with complex multiplication by $\OO_K$.

 From Abramovich's bound, since $[\PSL_2(\mathbb Z) :\Gamma_0(n)]\geq n+1$, it follows that $\gon_{\Q} X_0(n)\geq \gon_{\C} X_0(n)>\frac{325n}{2^{15}}$.

By \cite[Corollary 4.2.]{Kwon} it follows that $E_\Delta$ has an $n$-isogeny over $\Q(j(E_\Delta))$, a number field of degree $[\Q(j(E_\Delta)):\Q]=h_\Delta$.

By the arguments above and Minkowski's bound we have 
\begin{equation}\label{ineq}
d=h_{\Delta}\leq  \frac{2}{\pi}\sqrt{\Delta}     
 \leq \frac{4\sqrt{n}}{\pi}<\frac{1}{2}\cdot\frac{325n}{2^{15}} \leq \frac{1}{2}\cdot \gon_\Q X_0(n)
 \end{equation}
 as soon as 
\begin{equation}\label{ineq2}
\sqrt{n}>\frac{2^{18}}{325\pi},
\end{equation}
 i.e., when 
 $$n\geq 65920.$$
This (and hence \eqref{ineq}) is also satisfied as soon as 
$$\sqrt {n}\geq \frac{\pi d}{4}>256,$$
i.e., $d>201$.
 
 
For $d\leq 201$, it is enough to explicitly find an imaginary quadratic field $\Q(\sqrt{-n})$ with class number $d$ and satisfying
$$\frac{\Delta}{4}\geq n>\frac{2^{16}d}{325}$$
to prove the existence of a sporadic point of degree $d$ on $X_0(n)$. We verify this for all $d\leq 201$.

Furthermore, we verify the assumption that there exists a quadratic imaginary field with class number $d$ for all $d\leq 2166$ by searching through the LMFDB. \end{proof}

It is widely believed that the assumptions of \Cref{sporadic} are satisfied, i.e., that every positive integer is a class number. In particular, if one defines $\mathcal{F}(d)$ to be the number of imaginary quadratic fields with class number $d$, then it is expected (see \cite[(1.4)]{Soundarajan}, and also \cite {HJKMP2019}) that
$$\mathcal{F}(d)\asymp \frac{d}{\log d}.$$ Finally, we note that if one assumes the Generalized Riemann Hypothesis, then the assumptions of \Cref{sporadic} are satisfied for all $d\leq 10^6$ \cite[Section 9]{HJKMP2019}.

\section{A moduli problem with 2 torsion}

Let $R := \Spec \Z[1/2]$ and consider the category  $\Ell_R$ whose objects are pairs $(E,T)$ where $T$ is an $R$ scheme and $E$ is an elliptic curve over $T$; see \cite[Section 4.13]{KatzMazur85}. A morphism in $\Ell_R$ between $f_1: E_1 \to T_1$ and $f_2: E_2 \to T_2$ is a Cartesian square
	\[
	\begin{tikzcd}
	E_1 \arrow[r,"h"] \arrow[d,"f_1"] & E_2 \arrow[d, "f_2"] \\
	T_1 \arrow[r, "g"'] & T_2.
	\end{tikzcd}
	\] A moduli problem on $\Ell_R$ is a contravariant functor $\mathcal P : \Ell_R^{op} \to Sets$. We define $[\Gamma_1(2,\modcrvgt 2)] : \Ell/R \to Sets$ to be the moduli problem which sends $(E,T)$ to the set of pairs $P,Q \in E(T)$ such that $P$ is of exact order $2$ in all fibers and $Q$ is of exact order $>2$ in all fibers.

Define 
\begin{align*}
\Delta(b,c) &:= 16 b^2 (b + c - 1)^2(4b^2 + 4bc - 4b + c^2) \in \Z[b,c], \\
Y_1(2,\modcrvgt 2) &:= \Spec \Z\left[\frac 1 2, b,c, \frac 1 {\Delta(b,c)}\right], \\
E_1(2,\modcrvgt 2) &: y^2 =x^3 + c x^2 + (1-b-c)bx,  \\
P_1(2,\modcrvgt 2) &:= (0,0),\\
Q_1(2,\modcrvgt 2) &:= (b,b).
\end{align*}


\begin{theorem}
The curve $E_1(2,\modcrvgt 2)$ is an elliptic curve over $Y_1(2,\modcrvgt 2)$. The points $P_1(2,\modcrvgt 2)$ and $Q_1(2,\modcrvgt 2)$ on $E_1(2,\modcrvgt 2)$ have order $2$ respectively $>2$ in all fibers.
Furthermore, the triple $$(E_1(2,\modcrvgt 2),P_1(2,\modcrvgt 2),Q_1(2,\modcrvgt 2))$$ over $Y_1(2,\modcrvgt 2)$ is universal in the sense that it represents the moduli problem $[\Gamma_1(2,\modcrvgt 2)]$.
\end{theorem}

\begin{proof}
We start with $E/S$ written in Weierstrass form
\begin{equation}
    y^2+a_1xy+a_3y=x^3+a_2x^2+a_4x+a_6,
\end{equation}
together with two points $P,Q\in E(S)$ of order $2$ in all fibers and order $>2$ in all fibers, respectively. This equation can be mapped via an isomorphism $f$ to another equation in 
Weierstrass form, where the isomorphism $f$ is represented by a quadruple $(u,r,s,t)$, such that $f$ maps $(x,y)$ to 
$$((x-r)u^2, \; (y - s(x-r) - t)u^3),$$
see \cite[Proposition VII.1.3]{silverman}. The effect of this isomorphism on the coefficients $a_i$ is
\begin{equation} \label{eq:trans}
    \begin{split}a_1' &= (a_1 + 2s) u, \\
a_2' &= (a_2 - a_1s + 3r - s^2)u^2, \\
a_3' &= (a_3 + a_1r + 2t)u^3, \\
a_4' &= (a_4 + 2a_2r - a_1(rs+t) - a_3s + 3r^2 - 2st)u^4, \\
a_6' &= (a_6 - a_1rt + a_2r^2 - a_3t + a_4r + r^3 - t^2)u^6 .\end{split}
\end{equation}




We move the point $P$ to $(0,0)$ by selecting an appropriate $r$ (to move $x(P)$ to $0$) and $t$ (to move $y(P)$ to $0$, after $x$ has been moved to $0$). Now our equation is of the form
$$y^2+a_1'xy=x^2+a_2'x^2+a_4'x.$$ Notice that $a_6'=0$ since $(0,0)$ is a point on the curve and $a_3'=0$ since $(0,0)$ is of order 2.  We can move $a_1'$ to $0$ by selecting an appropriate $s$ in the quadruple corresponding to the isomorphism by \eqref{eq:trans}. Now the point $Q$ is of the form $(v,w)$ for some $v,w \in \OO_S(S)$ since it is of order $>2$ in all fibers. We can scale the coordinates of $Q$ by $(v,w)\mapsto(vu^2, wu^3)$. Hence, we can choose $u=w/v$ and letting $b:=v^3/w^2$, we get that $Q$ is now of the form $(b,b)$. 
Note here that $w,v,u\in \OO_S(S)^\times$, since otherwise $Q$ would be a point of order $2$ in the fiber where one of $w,v,u$ vanishes. 
Let 
$$y^2=x^3+a_2''x^2+a_4''x$$
be the equation we get after this transformation.
Let $c=a_2''$. Since $Q=(b,b)$ lies on this equation we get that
$$b^2=b^3+cb^2+ba_4'',$$
so $a_4''=\frac{b^2-b^3-cb^2}{b}=(1-b-c)b$, as claimed. 
\end{proof}

\begin{remark}
Jain gives in \cite{Jain10} a (two-variable) family of elliptic curves that is birational to $E_1(2,_>2).$. His family is given in terms of variables $t$ and $q$, and to get from $E_1(2,_>2)$ (given in terms of $b$ and $c$) to his family and vice versa, one uses the following transformations:
\begin{align*}
t := (-2b - c + 1)/(c-1);\\
q := b+1/2c+1/2-2b/(2b + c - 1);\\
b := (t+1)(qt+1)/t^2;\\
c := (t^2-2qt-2)/t^2.
\end{align*}
\end{remark}



\listoftodos[List of todo's]{}
\todo[inline]{remove todolist}

\begin{comment}
\todo[inline]{compare with Points of low height on elliptic surfaces with torsion | LMS Journal of Computation and Mathematics | Cambridge Core - \url{https://bit.ly/3C0gSCD} doi:10.1112/S1461157009000151
\begin{align*}
t := (-2*b - c + 1)/(c-1);\\
q := b+1/2*c+1/2-2*b/(2*b + c - 1);\\
b := (t+1)*(q*t+1)/t^2;\\
c := (t^2-2*q*t-2)/t^2;
\end{align*}

} 



\section{Overview}





Here is the list of the all solved cases:

\begin{itemize}
    \item Rank 0: 25,26,27,28,30,32,33,34,35,36,39,42,55,119,2x22,3x15
    \item Positive rank: 63, 65, 77,85,91,121,143,169,187,221,289
\end{itemize}
Number of cases solved: 31;\\
Number of cases to go: 2;\\

\bigskip

Solved but not yet included in this file: 
\begin{itemize}
    \item Maartens unpublished work: 30, 32, 34, 36
\end{itemize}


Remaining: 
\begin{itemize}
    %\item 63 ($\rk(J_1)=1, \rk(J_0)=0$, $g=5$).
    %\item 65 ($\rk(J_1)=\rk(J_0)=1$, $g=5$)
    \item $2\times 20$ (genus 9), $2\times 24$ (genus 13) (rank 0)
\end{itemize}


\subsection{Comments}
\begin{enumerate}
    \item $\gon_\Q X_0\leq 4$ for all remaining values.
    \item Using the Derickx-Stoll paper: we could eliminate the the case $n=77$ if $\gon_\Q X_1(77)> 32$, and $n=45,49, 51$ if $\gon_\Q X_1(n)> 16$ (possibly even $12$ if one can use the $T_2$ operator).
    \item Abramovich's bound gives the following pairs (45,9), (49,12), (51,12), (65,20), (77,29), (85,35), (91,40):
\end{enumerate}






\section{Useful stuff}
\todo[inline]{delete this section}
\subsection{Analytic rank of \texorpdfstring{$J_1/J_0$}{J1/J0}}

Check where the analytic rank of $J_0$ and $J_1$ differ:
\url{https://www.lmfdb.org/ModularForm/GL2/Q/holomorphic/?weight=2&char_order=2-&analytic_rank=1-&showcol=char_order.analytic_rank}

\subsection{Non-injectivity of the torsion}

The code:
\begin{verbatim}
from mdsage.maartens_sage_functions import *
G = Gamma11(3,5)
U,C = modular_unit_lattice(G, return_ambient=True, ambient_degree_zero=True)
C/U
\end{verbatim}

Shows that $J_1(3,15)(\Q(\zeta_{15}))[2^\infty]=\Z/2\Z \oplus (\Z/4\Z)^2 \oplus \Z/16\Z \oplus (\Z/32\Z)^2$
\end{comment}


\bibliographystyle{siam}
\bibliography{bibliography1}
\end{document}

